%%%%%%%%%%%%%%%%%%%%%%% Boundary-Layer Meteorology 2019 Template %%%%%%%%%%%%%%%%%%%%%%%%%
\begin{filecontents*}{example.pdf}
gsave
newpath
  20 20 moveto
  20 220 lineto
  220 220 lineto
  220 20 lineto
closepath
2 setlinewidth
gsave
  .4 setgray fill
grestore
stroke
grestore
\end{filecontents*}

\RequirePackage{fix-cm}
\documentclass[smallextended]{svjour3}       
\smartqed  
\journalname{Boundary-Layer Meteorology}
\usepackage{appendix}
\usepackage{amsmath}
\usepackage{graphicx}
\usepackage{lineno}
\usepackage{array}
\usepackage{longtable}
\usepackage{amssymb}
\usepackage{natbib}
\usepackage{hyperref}
\usepackage{xcolor}
\usepackage{tikz}
\usepackage{subcaption}
\usepackage{geometry}
\usepackage{booktabs,array}
\usepackage{placeins}
\usepackage{bm}
\usepackage{rotating}
\usepackage[table]{xcolor}

\hypersetup{
    colorlinks=true,
    linkcolor=blue,
    citecolor=black,
    urlcolor=blue
}

\geometry{
  top=2cm,
  bottom=2cm,
  left=2.5cm,
  right=2.5cm,
  marginparwidth=1.75cm
}
\usetikzlibrary{positioning, shapes, arrows.meta}

\setcitestyle{aysep={}}
\linenumbers

\newcommand*\patchAmsMathEnvironmentForLineno[1]{%
\expandafter\let\csname old#1\expandafter\endcsname\csname #1\endcsname
\expandafter\let\csname oldend#1\expandafter\endcsname\csname end#1\endcsname
\renewenvironment{#1}%
{\linenomath\csname old#1\endcsname}%
{\csname oldend#1\endcsname\endlinenomath}}% 
\newcommand*\patchBothAmsMathEnvironmentsForLineno[1]{%
\patchAmsMathEnvironmentForLineno{#1}%
\patchAmsMathEnvironmentForLineno{#1*}}%
\AtBeginDocument{%
\patchBothAmsMathEnvironmentsForLineno{equation}%
\patchBothAmsMathEnvironmentsForLineno{align}%
\patchBothAmsMathEnvironmentsForLineno{flalign}%
\patchBothAmsMathEnvironmentsForLineno{alignat}%
\patchBothAmsMathEnvironmentsForLineno{gather}%
\patchBothAmsMathEnvironmentsForLineno{multline}%
}

\begin{document}

\title{Development of Profile Assimilation Methods for Data-Driven Large Eddy Simulations}

\author{Arjun Ajay$^1$ \and
        Jagdeep Singh$^1$ \and
        Sebastiano Stipa$^2$ \and
        Philippe Beaucage$^3$ \and
        Joshua Brinkerhoff$^1$}

\institute{
    $^1$ The University of British Columbia---Okanagan, Canada \\
    \email{joshua.brinkerhoff@ubc.ca} 
    \and
    $^2$ von K\'arm\'an Institute for Fluid Dynamics, Sint-Genesius-Rode, Belgium
    \and
    $^3$ UL Renewables, USA 
}

\date{Received: DD Month YEAR / Accepted: DD Month YEAR}

\maketitle
%%%%%%%%%%%%%%%%%%%%%%%%%%%%%%%%%%%%%%%%%%%%%%%%%%%%
%%%                   Abstract                   %%%
%%%%%%%%%%%%%%%%%%%%%%%%%%%%%%%%%%%%%%%%%%%%%%%%%%%%

\begin{abstract}
Mesoscale-to-microscale coupling (MMC) extends the capability of large eddy simulations by introducing spatially and temporally varying mesoscale conditions in order to simultaneously model the large-scale atmospheric dynamics and the small-scale turbulent processes. This coupling is important for a range of applications, including wind energy forecasting, wildfire spread prediction, and urban climate modeling. In this study, we introduce two new offline MMC methods within the class of profile assimilation techniques. The first method, wavelet-based profile assimilation (WPA), applies multi-resolution decomposition to compute the internal forcing for the momentum and temperature equations, while the second method is a hybrid approach that combines a physics-based geostrophic balance for momentum source terms with wavelet-based profile assimilation for temperature source terms. Unlike previous profile assimilation methods, such as direct and indirect profile assimilation (DPA, IPA), which apply error-based corrections to both momentum and temperature equations, the proposed hybrid method avoids momentum error forcing, which simplifies the MMC algorithm and implementation. To demonstrate our approach, a composite mesoscale dataset derived from the American WAKE experimeNt (AWAKEN) is used to add mesoscale forcing to large eddy simulations conducted over flat terrain for a full diurnal cycle. The composite dataset combines experimental observations from scanning lidar measurements spanning heights from $100$ to $1100~\mathrm{m}$, virtual tower data, and mesoscale output from Weather Research and Forecasting (WRF) model simulations. Our results show that the new assimilation strategies yield mean flow profiles and turbulence statistics that achieve substantially better accuracy relative to previous profile assimilation techniques. The proposed framework offers a flexible and physically consistent method for coupling mesoscale information into large eddy simulations across a range of atmospheric boundary layer conditions.

\keywords{Data assimilation, Diurnal cycle, Large eddy simulation,  Mesoscale-to-microscale coupling, American WAKE experimeNt }
\end{abstract}

%%%%%%%%%%%%%%%%%%%%%%%%%%%%%%%%%%%%%%%%%%%%%%%%%%%%
%%%                   Intro                      %%%
%%%%%%%%%%%%%%%%%%%%%%%%%%%%%%%%%%%%%%%%%%%%%%%%%%%%
\section{Introduction}
\label{sec:intro}
Microscale models, such as large-eddy simulations (LES), are among the most prominent numerical tools for investigating turbulent fluxes of mass, momentum, heat, and pollutants in the atmospheric boundary layer (ABL) \citep{stoll2020large}. Operating at meter-scale resolutions, these models can resolve fine-scale flow features, including boundary-layer processes, turbulence characteristics, and terrain-induced flow patterns. However, their high computational cost restricts their domain size to at most $\mathcal{O}$(10 km), preventing adequate representation of larger-scale atmospheric dynamics. Consequently, LES studies struggle to capture crucial mesoscale meteorological phenomena that drive microscale flows, such as frontal passages, thunderstorm outflows, diurnal land-sea breezes, and low-level jets (LLJs). Additionally, most LES solvers lack appropriate models to represent mesoscale weather processes, including radiative transfer, cloud microphysics, and land-surface interactions \citep{haupt2023lessons}. As a result, a majority of the LES studies of ABLs have focused on idealized, canonical flow scenarios. A recent review by \citet{stoll2020large} provides a comprehensive overview of such canonical cases and details the 50-year development of LES in studying boundary layer processes. Efforts are ongoing to incorporate non-stationary, large-scale mesoscale weather phenomena into LES \citep{Haupt2019, haupt2023lessons, sanz2017mesoscale}. However, combining these fine-scale flow dynamics with broader atmospheric processes remains a central challenge in achieving realistic simulations of earth's atmospheric boundary layer. 

By comparison, General Circulation Models (GCMs), with horizontal resolutions ranging from tens to hundreds of kilometers, effectively capture large-scale weather and climate patterns but fail to resolve finer turbulent structures that are important for boundary-layer dynamics. Mesoscale models, including the Weather Research and Forecasting (WRF) model \citep{michalakes2001, michalakes2005weather} and the Energy Research and Forecasting (ERF) model \citep{almgren2023erf}, function at kilometer-scale resolutions and partially address this resolution disparity. However, they still lack the ability to fully resolve turbulence and subgrid-scale flow interactions. 

Mesoscale-to-microscale coupling (MMC) is a modeling approach that aims to simultaneously model mesoscale meteorological phenomena and the small-scale turbulent processes by enforcing realistic mesoscale conditions within high-resolution, microscale simulations. The purpose of MMC is to drive the microscale model with information obtained from mesoscale models, thereby capturing the full range of atmospheric motions and improving the understanding of boundary-layer turbulence and its interaction with larger-scale weather patterns. Applications of MMC span a wide range of atmospheric flows, including wind energy \citep{sanz2017mesoscale, sanz2017methodology, quon2024measurement}, wildfire dynamics \citep{nakayama2015large}, urban meteorology \citep{baklanov2009multi, liu2012study}, and pollutant dispersion \citep{nakayama2015large}.

Although various strategies exist for MMC \citep{haupt2023lessons}, this study adopts an offline coupling approach where mesoscale and microscale solvers operate independently. Within this framework, mesoscale information can be integrated into the microscale model via boundary coupling or using internal source terms. In the boundary coupled approach, mesoscale model outputs provide time-dependent inflow conditions at the boundaries of the microscale domain. A key challenge in this method is the transition of momentum, heat, and turbulence across scales. Because mesoscale models do not resolve fine-scale turbulence, allowing the natural energy cascade to generate these smaller structures often requires large domains or significant spin-up periods \citep{meneveau2019big}. This issue can be addressed by introducing synthetic turbulence at the inflow to accelerate the formation of small-scale structures \citep{mirocha2014resolved, munoz2015stochastic,munoz2018generation,quon2018enrichment}. Alternatively, internal forcing can be applied by incorporating mesoscale forcing as source terms in the microscale governing equations. In this case, synthetic turbulence at the inflow can be avoided by using a precursor simulation \citep{stevens2014concurrent, stipa2023tosca} to generate realistic turbulence. This method is commonly employed in canonical ABL simulations together with periodic boundary conditions in the horizontal directions, thereby assuming horizontal homogeneity, which is a reasonable approximation under simple orographic and meteorological conditions.

Two prominent methods for applying internal forcing in microscale simulations have been proposed in the literature \citep{allaerts2020development}. The first one is the mesoscale budget component approach which extracts momentum and temperature transport budget equations from a mesoscale solver, such as numerical weather prediction (NWP) models. These budget terms include the background horizontal pressure gradient and large-scale advection. The extracted terms are then imposed directly into the microscale governing equations. Several studies have used this approach \citep{baas2010design,neggers2012continuous,schalkwijk2015year,heinze2017evaluation,sanz2017methodology,olsen2018mesoscale}. However, an evaluation by \citet{draxl2021coupling} highlighted challenges with this approach, especially regarding the small-scale spatio-temporal variability of the mesoscale transport budget terms. Due to the disparate resolutions of the mesoscale and microscale solutions, it is not possible to distinguish whether these fluctuations represent real physical processes or are numerical artifacts. Non-physical fluctuations require filtering to prevent contamination of the microscale solution. The approach also encounters model drift, where the microscale solution deviates from mesoscale trends, often requiring non-physical nudging to ensure consistency.

An alternative approach, known as the profile assimilation method, was developed by \citet{allaerts2020development} and has been successfully used in several subsequent studies \citep{wrba2024,allaerts2023using,quon2024measurement}. This method utilizes primitive mesoscale variables (e.g., wind speed and virtual potential temperature) rather than relying on transport-equation budget components. Similar to data assimilation, this method provides the microscale solver with the time-height history of primitive flow variables, enabling it to estimate internal forcing terms dynamically. In this framework, \citet{allaerts2020development} proposes two profile assimilation techniques for coupling mesoscale inputs with microscale solvers: direct profile assimilation (DPA) and indirect profile assimilation (IPA).

DPA applies a forcing proportional to the difference between the plane-averaged primary quantities $\langle \varphi_{\text{\tiny les}} \rangle$ of the microscale solver and either standard mesoscale solver output or observational data $\varphi_{\text{\tiny m}}$ \citep{allaerts2020development,draxl2021coupling}. Previous studies have noted an artificial buildup of turbulent kinetic energy (TKE) when employing DPA \citep{quon2024measurement,allaerts2023using,allaerts2020development}, which arises from the over-constrained nature of DPA, given that the solver has limited capacity to adjust the mean flow via turbulent mixing. 

On the other hand, IPA fits a polynomial function to the DPA error profile. The polynomial fitting acts as a low-pass filter that smooths the corrections in regions where the mean microscale profile deviates from observational data, allowing turbulent kinetic energy (TKE) to remain closer to physically realistic levels. Nonetheless, IPA is prone to over-smoothing of critical flow variations due to the order of the polynomial function used \citep{allaerts2020development}. 

Addressing the deficiencies of DPA and IPA is essential to improve how boundary-layer turbulence is captured and to improve the overall reliability of mesoscale-to-microscale coupling across different atmospheric conditions. In this study, we developed two new internal forcing approaches to address the challenges identified in previous approaches. The first method, called wavelet-based profile assimilation (WPA), applies multi-resolution analysis (MRA) \citep{mallat1989theory,debnath2017construction,singh2024impact} to the direct error profile between mesoscale and microscale quantities. Acting as a scale-aware filter, it preserves vertical structure and localized turbulent features while avoiding over-smoothing of the source terms. Temperature forcing is applied using the wavelet-based approach to the virtual potential temperature field. The second method, hybrid geostrophic-wavelet profile assimilation (HGWPA), computes mesoscale momentum forcing from the large-scale pressure gradient based on geostrophic balance, with the geostrophic wind varying with height under baroclinic conditions. Unlike error-based techniques, HGWPA emphasizes physics-driven forcing to maintain mesoscale-microscale coherence. Both WPA and HGWPA methods rely on primary mesoscale variables derived from standard mesoscale model data or field measurement data and avoid the use of budget terms from mesoscale transport equations, which we show alleviates the challenges inherent in existing profile assimilation techniques.

The remainder of this paper is organized as follows. Sect.~\ref{sec:methods} describes the methods, including the background of existing MMC algorithms and the proposed wavelet-based and hybrid profile assimilation approaches. Sect.~\ref{sec:computationalModel} outlines the observational data used to drive the mesoscale dynamics and explains the setup of the large eddy simulations used to demonstrate our MMC approach. Sect.~\ref{sec:results} presents the results of these simulations, evaluating how the proposed methods capture the mean profiles, turbulent kinetic energy, and mesoscale forcing consistency in an ABL over flat terrain. Sect.~\ref{sec:discussion} compares the performance of the profile assimilation methods, ranks them using error metrics, and discusses the turbulent characteristics based on energy spectra. Finally, Sect.~\ref{sec:conclusions} summarizes the main findings and suggests directions for future work.
%%%%%%%%%%%%%%%%%%%%%%%%%%%%%%%%%%%%%%%%%%%%%%%%%%%%
%%%                   METHOD                     %%%
%%%%%%%%%%%%%%%%%%%%%%%%%%%%%%%%%%%%%%%%%%%%%%%%%%%%
\section{Meso-to-micro scale coupling methodology}
\label{sec:methods}
\subsection{\textbf{Background}}
\label{subsec:simnum}
We employ a multiscale MMC model within the TOSCA LES code \citep{stipa2023tosca} to capture the energetic turbulent motions within a dynamically evolving ABL. Our methodology is informed by the framework proposed by \cite{allaerts2020development} for deriving the mesoscale source terms. The momentum and virtual potential temperature governing equations expressed in Cartesian co-ordinates are: 
%
% Momentum Equation (split over two lines)
\begin{equation}
\begin{split}
\frac{\partial \overline{u}_i}{\partial t} + \frac{\partial}{\partial x_j}\bigl(\overline{u}_j \overline{u}_i\bigr)
&= -\frac{1}{\rho_0}\frac{\partial \overline{p}^*}{\partial x_i} 
- \frac{\partial \tau_{ij}^{D}}{\partial x_{j}} \\
&\quad  + \left[1 - \left(\frac{\overline{\theta}-\theta_0}{\theta_0}\right)\right]g_i - 2\epsilon_{ijk}\Omega_j \overline{u}_k + F_{u_i},
\end{split}
\label{eq:momentumCartesian}
\end{equation}
% Temperature Equation
\begin{equation}
    \frac{\partial \overline{\theta}}{\partial t} + \frac{\partial}{\partial x_j}\bigl(\overline{u}_j \overline{\theta}\bigr)
= \frac{\partial \tau_{\theta j}}{\partial x_j} + F_{\theta} ,
\label{eq:temperatureCartesian}
\end{equation}
%
where, the ($\overline{\cdot}$) represents spatial filtering using a filter of the LES grid scale $\Delta_{\hbox{\tiny les}}$ such that $\overline{u}_i$ is the filtered Cartesian velocity and $\overline{\theta}$ is the filtered virtual potential temperature. $\tau_{i j}^D=\overline{u_i u_j}-\overline{u}_i \overline{u}_j$ is the deviatoric part of the subgrid scale (SGS) stress tensor and $\tau_{\theta j}=\overline{u_j \theta}-\overline{u}_i \overline{\theta}$ is the subgrid heat flux. The quantity $\overline{p}^*/\rho_0$, where $\rho_0$ is the density of dry air, is the filtered modified kinematic pressure with the trace of the SGS stress tensor absorbed into it. $\Omega_j$ is the Earth's angular velocity vector and $\varepsilon_{ijk}$ is the Levi-Civita permutation tensor. The buoyancy term from the Boussinesq approximation comprises $\overline{\theta}$, $\theta_0$, which is the reference virtual potential temperature, and $g_i$ is the gravitational acceleration vector. The additional terms $F_{u_i}$ and $F_{\theta}$ in Eqs.~\ref{eq:momentumCartesian} and \ref{eq:temperatureCartesian}, respectively, represent the mesoscale forcing sources that drive the mean microscale velocity and virtual potential temperature fields to their mesoscale mean values. 

In ABL simulations, microscale solvers commonly employ anisotropic grids with finer resolution in the vertical direction. However, most eddy-viscosity-based turbulence models assume isotropic grid spacing. Notable exceptions, such as the anisotropic minimum dissipation model (AMD) model \citep{Moser2021}, explicitly account for grid anisotropy, improving their accuracy in such configurations. In this study, the subgrid terms $\tau_{i j}$ and $\tau_{\theta j}$ are modeled using AMD \citep{rozema2015, AbkarAMD2016}.  The subgrid tensor $\tau_{ij}^{D}$ is modeled as
\begin{equation}
\tau_{i j}^D=\tau_{i j}-\frac{1}{3} \delta_{i j} \tau_{k k}=-2 \nu_tS_{i j},
\end{equation}
where $\nu_{t}$ is the subgrid scale eddy viscosity and $S_{ij}=\frac{1}{2}$ $\left(\partial_i \bar{u}_j+\partial_j \bar{u}_i\right)$ is the resolved rate-of-strain tensor. The subgrid heat flux  $\tau_{\theta j} $ is modeled as $\tau_{\theta j}=-\kappa_{t} \partial_j \bar{\theta}$ where $\kappa_{t}$ is the subgrid scale eddy diffusivity. The AMD model is adapted for stratified atmospheres \citep{Abkar2017} with scaled velocity and displacement vectors to ensure the eddy viscosity is frame invariant in anisotropic grids \citep{Verstappen2018}. It provides $\nu_t$ as 
\begin{equation}
\nu_t=-\left(C \Delta^2\right)\frac{\left(\hat{\partial}_k \hat{u}_i\right)\left(\hat{\partial}_k \hat{u}_j\right) \hat{S}_{i j}+\hat{\delta}_{i 3} \beta\left(\hat{\partial}_k \hat{u}_i\right) \hat{\partial}_k(\bar{\theta})}{\left(\hat{\partial}_l \hat{u}_m\right)\left(\hat{\partial}_l \hat{u}_m\right)},
\label{eq:subgridviscosity}
\end{equation}
%
where $C$ is the modified Poincar\'e constant. The scaled spatial coordinate, velocity, and velocity gradient are given by 
$\hat{x}_i=\frac{x_i}{\Delta_i}, \hat{u}_i=\frac{\overline u_i}{\Delta_i}, \hat{\partial}_i \hat{u}_j=\frac{\Delta_i}{\Delta_j} \partial_i \overline u_j$, respectively, while $\hat{\delta}_{i 3}=\frac{\delta_{i 3}}{\Delta_3}$ is the scaled Kronecker delta tensor and $\hat{S}_{i j}=\frac{1}{2}\left(\hat{\partial}_i \hat{u}_j(\hat{x}, t)+\hat{\partial}_j \hat{u}_i(\hat{x}, t)\right)$ is the scaled rate of strain tensor. Here $\Delta_i$ is the filter width and $\beta$ is the buoyancy parameter given by $\beta = \frac{g}{\theta_0}$. Repeated indices in Eq.~(\ref{eq:subgridviscosity}) indicate summation over the spatial dimensions (Einstein summation convention), so that all tensor contractions are fully summed and the resulting eddy viscosity $\nu_t$ is a scalar quantity.

Eddy diffusivity $\kappa_t$ is given by 
\begin{equation}
\kappa_t=\left(C \Delta^2\right)\frac{-\left(\hat{\partial}_k \hat{u}_i\right)\left(\hat{\partial}_k \overline \theta\right) \hat{\partial_i} \overline \theta}{\left(\hat{\partial}_l \overline \theta\right)\left(\hat{\partial}_l \overline \theta\right)},
\label{eq:subgriddiffusivity}
\end{equation}
In both Eqs.~\ref{eq:subgridviscosity} and \ref{eq:subgriddiffusivity}, mean filter width $\Delta$ is computed as the square root of the harmonic mean of the square of the filter width in each direction:
\begin{equation}
\frac{1}{\Delta^2}=\frac{1}{3}\left(\frac{1}{\Delta_x^2}+\frac{1}{\Delta_y^2}+\frac{1}{\Delta_z^2}\right).
\end{equation}
The filter widths ($\Delta_x, \Delta_y, \Delta_z$) are chosen as the local grid size in the Cartesian coordinate directions, while the modified Poincar\'e constant is taken as 1/3 \citep{Abkar2017}.

Finally, to complete the modeling framework, MMC is achieved using a profile assimilation technique~\citep{allaerts2020development}, which incorporates time- and height-dependent velocity and virtual potential temperature data to compute mesoscale forcing terms. The following sections briefly review some of the prominent profile assimilation methodologies (Sect. \ref{sec:exitingPA}), and introduce the new assimilation techniques (Sect. \ref{sec:proposedPA}) within this framework.

\subsection{\textbf{Existing profile assimilation methods}}
\label{sec:exitingPA}
\subsubsection{Direct profile assimilation}
\label{sec:DPA}
Direct profile assimilation approach estimates the forcing terms, $F_{u_i}$ and $F_{\theta}$ for Eqs.~\ref{eq:momentumCartesian} and~\ref{eq:temperatureCartesian} using an error term that is imposed via a proportional controller. The error term is computed as the difference between the plane-averaged primary variables $\langle \bar\varphi_{\hbox{\tiny LES}}\rangle$, and the mesoscale data $\varphi_{\hbox{\tiny m}}$ at each vertical level such that the forcing between the different vertical levels remains uncorrelated. This error term is given by:
\begin{equation}
e_{\varphi}(z, t) = \bigl(\varphi_{\hbox{\tiny m}}(z,t) - \langle\bar\varphi_{\hbox{\tiny LES}}\rangle(z,t)\bigr).
\label{eq:errorterm}
\end{equation}
The forcing terms $F_{u_i}$ and $F_{\theta}$ are then computed as:
\begin{equation}
\mathcal{F}_{\varphi}(z,t) = \mathcal{K}_p e_{\varphi}(z,t),
\label{dpaForcing}
\end{equation}
where $\varphi \in \{u, v, \theta \}$, and $\mathcal{K}_p = 0.2~\mathrm{s^{-1}}$ is the proportional controller gain used consistently in all simulations presented in this article. Sensitivity tests varying $\mathcal{K}_p$ showed negligible impact on the results, consistent with prior studies \citep{allaerts2020development}.

\subsubsection{Indirect profile assimilation}
\label{subsec:ipa}
An alternative to the DPA approach is the indirect profile assimilation methodology \citep{allaerts2020development} that introduces a vertical smoothing of the error profile $e_\varphi(z,t)$. Instead of applying the instantaneous error at each vertical level, the indirect technique fits a polynomial of order \(m\) to the error $e_\varphi(z,t)$ (Eq.~\ref{eq:errorterm}) at every timestep, thus correlating the corrections across the entire vertical length. Formally, \(e_{\varphi}(z,t)\) is approximated by a polynomial function
\(\widetilde{e}_{\varphi}(z,t) = Z\,\hat{\boldsymbol{\beta}}_{\varphi}(t)\), where \(\hat{\boldsymbol{\beta}}_{\varphi}(t)\)
is a vector of polynomial coefficients and \(Z\) is the matrix formed by evaluating the basis \( \{\,1,\,z,\,z^{2},\dots,z^{m}\} \) at the vertical grid points \(\{z_{1},z_{2},\dots,z_{n}\}\), namely,
%
\begin{equation*}
  Z \;=\;
  \begin{bmatrix}
    1      & z_{1}   & z_{1}^{2} & \dots  & z_{1}^{m}\\
    1      & z_{2}   & z_{2}^{2} & \dots  & z_{2}^{m}\\
    \vdots & \vdots  & \vdots    & \ddots & \vdots   \\
    1      & z_{n}   & z_{n}^{2} & \dots  & z_{n}^{m}
  \end{bmatrix}.
\end{equation*}
%
At each time step \(t\), the polynomial coefficients \(\hat{\boldsymbol{\beta}}_{\varphi}(t)\) are found by
minimizing a weighted least-squares cost:
%
\begin{equation}
  \hat{\boldsymbol{\beta}}_{\varphi} \;=\;
  \underset{\boldsymbol{\beta}}{\mathrm{arg\,min}}\,
    \bigl\|\,W^{\!1/2}\,\bigl(e_{\varphi} - Z\,\boldsymbol{\beta}\bigr)\bigr\|_{2}^{2},
\end{equation}
%
where \(W\) is a diagonal matrix of positive weights \(\{w_{k}\}\).  The solution to this unconstrained quadratic
problem is computed explicitly by
%
\begin{equation}
  \hat{\boldsymbol{\beta}}_{\varphi}
  \;=\;
  \bigl(Z^{\mathsf{T}}\,W\,Z\bigr)^{-1}\,Z^{\mathsf{T}}\,W\,e_{\varphi}.
\end{equation}
%
Once the polynomial fit \(\widetilde{e}_{\varphi}\) is obtained, it is employed in the forcing term:
%
\begin{equation}
  F_{\varphi}(z,t) \;=\; K_{p}\,\widetilde{e}_{\varphi}(z,t).
  \label{ipaForcing}
\end{equation}
%
Because the assimilation correction \(\widetilde{e}_{\varphi}\) is derived from a smoothed (polynomial-fitted) version of the vertical error profile $e_\varphi$, the resulting forcing reflects the overall vertical trend rather than discrete pointwise mismatches. In this sense, IPA acts as a low-pass filter that enhances vertical spatial coherence. Typical choices for the polynomial order \(m\) range from linear to cubic or higher, providing increased flexibility in representing vertical gradients. The weight distribution \(\{w_{k}\}\) can also be tuned to prioritize certain height ranges such as within the ABL. Together, the polynomial order and weighting scheme serve as additional controller parameters in this methodology. In the present article, we employ a cubic polynomial order while using a uniform weighting scheme.  

In the following subsection, we present two new data assimilation approaches to compute the forcing $F_{u_i}$ and $F_{\theta}$ terms. The first method relies on multi-resolution analysis of the source term using wavelet transforms which accounts for vertical coherence as well as capturing the local variations and is discussed in Sect.~\ref{subsub:wpa}. Later, in Sect.~\ref{sec:geoDevelop}, an alternate method to the conventional profile assimilation techniques based on the geostrophic balance is introduced. 

%\subsection{New data assimilation methods}
\subsection{\textbf{Proposed meso-to-microscale approaches}}
\label{sec:proposedPA}
\subsubsection{Wavelet profile assimilation}
\label{subsub:wpa}
Wavelet profile assimilation uses wavelet-based multi-resolution analysis (MRA) \citep{mallat1989theory,debnath2017construction} to decompose the error \( e_\varphi(z,t) \) into scale-dependent components, enabling a reconstruction strategy that differs from the polynomial fitting approach. Consider an error term \( e_{\varphi}(z,t) \) as defined in Eq.~\ref{eq:errorterm}, where \( z \) denotes the vertical coordinate and \( t \) represents time. To analyze the scale-dependent structure of this error field in the vertical direction, a wavelet transform is applied at each time \( t \). The continuous wavelet transform (CWT) of \( e_{\varphi}(z,t) \) along the vertical direction \( z \) is defined as
\begin{equation}
\mathbb{W}_{a,b}(t) = \frac{1}{\sqrt{|a|}} \int_{-\infty}^{+\infty} e_{\varphi}(z,t) \psi \left( \frac{z - b}{a} \right) \mathrm{d}z,
\label{eq:wt}
\end{equation}
where \( \psi(z) \) is the mother wavelet function \citep{mallat1999wavelet}, and \( \mathbb{W}_{a,b}(t) \) denotes the wavelet coefficient at scale \( a \), location \( b \), and time \( t \). The parameters \( a \) and \( b \) correspond to the dilation (scale) and translation (location) of the wavelet, respectively. The family of scaled and shifted wavelets is given by
\begin{equation}
\psi_{a,b}(z) = \frac{1}{\sqrt{|a|}} \psi\left( \frac{z - b}{a} \right).
\end{equation}
These transformations enable multi-resolution analysis in both space and time, allowing localized features at different vertical scales to be identified and analyzed.

% Illustrating the Mexican hat wavelet
For illustration of the CWT, Fig.~\ref{fig:mexicanhat} presents a typical Mexican hat mother wavelet, which is the second derivative of a Gaussian function with effective compact support (within \([-z, z]\)), along with several of its scaled and shifted forms. The translation parameter \( b \) determines the spatial localization center, such that each \( \psi_{a,b}(z) \) is concentrated around \( z = b \). The scale parameter \( a \) governs the spatial frequency content; for \( |a| \ll 1 \), the wavelet is narrow and sensitive to fine-scale (high-frequency) features in \( e_{\varphi}(z,t) \); for \( |a| \gg 1 \), it becomes wider and captures coarse-scale (low-frequency) structures. The figure includes both fine and coarse scale representations, highlighting the multiscale and localized characteristics of wavelet functions.

\begin{figure}[htb!]
    \centering
    \includegraphics[width=0.8\textwidth]{Plots/waveletSchematic.pdf}
\caption{Typical Mexican-hat mother wavelet for the continuous wavelet transform (top panel) at scale \(a=1\) and translation \(b=0\). 
The bottom panel shows a family of wavelets obtained by varying the scale \(a\) and translation \(b\), demonstrating spatial localization from fine to coarse scales.}

    \label{fig:mexicanhat}
\end{figure}

% Introducing the DWT
In ABL simulations, the error field is available at discrete vertical levels \( k \). We apply the discrete wavelet transform (DWT) within a MRA framework \citep{mallat1989theory} to decomposes the signal into scale-dependent detail and approximation components, capturing high- and low-frequency features, respectively.  The detail coefficients, representing high-frequency components, are defined as
\begin{equation}
\mathbb{W_D}_{m,n}(t) = \sum_{k} e_{\varphi}(k,t) \psi_{m,n}(k),
\label{eq:DWT_detail}
\end{equation}
and the approximation coefficients, representing low-frequency components, as
\begin{equation}
\mathbb{W_A}_{m,n}(t) = \sum_{k} e_{\varphi}(k,t) \phi_{m,n}(k),
\label{eq:DWT_approx}
\end{equation}
where \( \mathbb{W_D}_{m,n}(t) \) and \( \mathbb{W_A}_{m,n}(t) \) are the detail and approximation coefficients at discrete scale level \( m \), translation index \( n \), and time \( t \). This forward transform (``analysis step'') projects the discrete error signal \( e_{\varphi}(k,t) \) onto the wavelet basis functions \( \psi_{m,n}(k) \) and scaling functions \( \phi_{m,n}(k) \), yielding the corresponding detail and approximation coefficients. Here, the wavelet basis functions are defined as:
\begin{equation}
\psi_{m,n}(k) = a_0^{-m/2} \psi(a_0^{-m} k - n b_0),
\label{eq:wavelet_basis}
\end{equation}
and the scaling functions as:
\begin{equation}
\phi_{m,n}(k) = a_0^{-m/2} \phi(a_0^{-m} k - n b_0),
\label{eq:scaling_basis}
\end{equation}
where \( \psi \) and \( \phi \) are the mother wavelet and scaling function, respectively, with standard DWT parameters \( a_0 = 2 \) and \( b_0 = 1 \) ensuring a dyadic grid for computational efficiency. This formulation results in a binary dilation of \( 2^{-m} \) and a translation of \( n 2^m \) \citep{debnath2017construction} such that the index \( m \) controls the scale (with higher \( m \) corresponding to coarser scales), while \( n \) determines the spatial localization of the wavelet.

% Inverse DWT and orthonormality
Appropriate choices of \( \psi \), \( \phi \), \( a_0 \), and \( b_0 \) ensure that \( \{\phi_{m,n}, \psi_{m,n} \mid m,n \in \mathbb{Z}\} \) form an orthonormal basis for the discrete signal space, enabling exact reconstruction via the inverse transform (``synthesis step''):
\begin{equation}
\hat{e}_{\varphi}(k,t) = \sum_{n} \mathbb{W_A}_{m_0,n}(t) \phi_{m_0,n}(k) + \sum_{m=0}^{m_0-1} \sum_{n} \mathbb{W_D}_{m,n}(t) \psi_{m,n}(k),
\label{eq:DWT_inverse}
\end{equation}
where \( \hat{e}_{\varphi}(k,t) \) is the reconstructed signal and \( m_0 \) is the total number of decomposition levels. The MRA framework structures this decomposition hierarchically, separating the signal into nested subspaces of increasing coarseness. Wavelet functions act as band-pass filters, capturing localized high-frequency details, while scaling functions capture low-frequency approximations. By selectively filtering these coefficients, MRA enables the isolation and reconstruction of errors at specific scales. Fig.~\ref{fig:mradwt} shows a representative schematic of an MRA with three levels of decomposition, where the first level splits the error signal into a coarse approximation level (\( \mathbb{W_A}_{1,n}(t) \)) and fine detail levels (\( \mathbb{W_D}_{1,n}(t) \)). Subsequent levels split the approximation signal at the previous level (\( \mathbb{W_A}_{m-1,n}(t) \)) to approximation (\( \mathbb{W_A}_{m,n}(t) \)) and detail levels (\( \mathbb{W_D}_{m,n}(t) \)), leading to the original signal represented as the coarsest approximation level (\( \mathbb{W_A}_{3,n}(t) \)) plus a series of detail levels (\( \mathbb{W_D}_{1-3,n}(t) \))  as shown in Eq. \ref{eq:DWT_inverse}.

\begin{figure}[htb!]
    \centering
    \includegraphics[width=0.5\linewidth]{waveletSchematic.png}
    \caption{A typical multi-resolution analysis (MRA) with three levels of decomposition. At each level, the error signal \( e_{\varphi}(k,t) \) is split into low-frequency approximation coefficients (\( \mathbb{W_A}_{m,n}(t) \)) and high-frequency detail coefficients (\( \mathbb{W_D}_{m,n}(t) \)).}
    \label{fig:mradwt}
\end{figure}

In this study, the Daubechies wavelet \citep{Daubechies1992} of order 7 (Db7) is used as the mother wavelet for its balance of smoothness and compact support, which is suitable for atmospheric profiles; this choice was validated through a sensitivity study of Daubechies wavelets from order 5 to 10 with varying decomposition levels (not shown here). A comprehensive exploration of other wavelet families, their orders, and corresponding decomposition levels is left for future work, which will employ machine learning-based hyper-parameter optimization to identify the most suitable wavelet configurations. Additionally, a symmetric signal extension is applied at the vertical domain boundaries to mitigate edge effects. This technique extends the signal by reflecting it about the boundary, preserving continuity and minimizing edge artifacts in the wavelet transform. These edge effects were previously observed to affect the forcing terms close to the surface for IPA \citep{allaerts2020development}.  

Applying MRA to the error profiles within WPA allows for flexible control over the level of detail in the reconstruction. A full reconstruction, including all detail coefficients, corresponds to a DPA-like formulation of the error term. Conversely, retaining only the coarsest approximation coefficients aligns with an IPA-like formulation. This adaptability is a key advantage of WPA, as it permits tailoring the error term formulation to account for factors such as boundary layer stability, boundary layer height, and atmospheric turbulence. In this study, we utilize six decomposition levels (\( m_0 = 6 \)) for Db7 wavelet such that the reconstructed error \( \hat{e}_{\varphi}(k,t) \) is formed using only the low-frequency (approximation) coefficients (\( \mathbb{W_A}_{m_0,n}(t) \)), thereby retaining only the essential large-scale structures in the vertical direction. Thus, the reconstructed error is
\begin{equation}
\hat{e}_{\varphi}(k,t) = \sum_n \mathbb{W_A}_{m_0,n}(t) \phi_{m_0,n}(k),
\end{equation}
and finally, the forcing term is given by
\begin{equation}
\mathcal{F}_{\varphi}(k,t) = \mathcal{K}_{p} \hat{e}_{\varphi}(k,t) = \mathcal{K}_{p} \sum_n \mathbb{W_A}_{m_0,n}(t) \phi_{m_0,n}(k),
\label{wpaForcing}
\end{equation}
where $\mathcal{K}_p = 0.2~\mathrm{s^{-1}}$ is a proportionality constant, chosen to be the same as in section \ref{sec:DPA}.
%%%%%%%%%%%%%%%%%%%%%%%%%%%%%%%%%%%%%%%%%%%%%%%%%%%%%%%%%%%%%%%%%%%%%%%%
\subsubsection{Hybrid geostrophic-wavelet profile assimilation}
\label{sec:geoDevelop}
HGWPA aims to estimate the internal mesoscale momentum forcing by computing the large-scale pressure gradient force, based on the principle of geostrophic balance. This balance represents an equilibrium between the Coriolis force and the horizontal pressure gradient force in the free-stream region above the ABL. Under steady-state conditions, geostrophic balance yields the large-scale pressure gradient as

\begin{equation}
    -\frac{\partial p_\infty}{\partial x} = -f_c\,G_y, 
    \quad 
    -\frac{\partial p_\infty}{\partial y} = f_c\,G_x,
    \label{eq:geostrophicBalanceSteady1}
\end{equation}
where \( f_c = 2\Omega \sin\phi \) is the Coriolis parameter, \( \Omega \) denoting the Earth's angular velocity, \( \phi \) is the geographic latitude, and the variables \( G_x \) and \( G_y \) represent the components of the geostrophic wind vector.

However, when the geostrophic wind varies in time, Eq. \ref{eq:geostrophicBalanceSteady1} no longer represents a strict balance but instead reflects a dynamical adjustment between the pressure gradient and Coriolis forces. This modifies Eq.~\eqref{eq:geostrophicBalanceSteady1} into its unsteady form as
\begin{equation}
    -\frac{\partial p_\infty(t)}{\partial x} = \frac{dG_x(t)}{dt} - f_c\,G_y(t), 
    \quad 
    -\frac{\partial p_\infty(t)}{\partial y} = \frac{dG_y(t)}{dt} + f_c\,G_x(t).
    \label{eq:geostrophicBalanceUnsteady2}
\end{equation}

Moreover, under realistic atmospheric conditions, horizontal temperature gradients can give rise to variations in geostrophic wind with height, a condition known as baroclinicity~\citep{sorbjan2004large,holton2013introduction,floors2015effect,momen2018modulation}, implying that the large-scale pressure gradient is no longer vertically uniform. Above the ABL, Eq. \ref{eq:geostrophicBalanceUnsteady2} remains valid on each horizontal plane, with the pressure gradient force balanced by the Coriolis force and inertial term. However, within the ABL, this balance is modified due to the influence of frictional drag arising from turbulent mixing. Consequently, the computation of the large-scale pressure gradient must include the divergence of turbulent momentum fluxes, leading to the following form:

\begin{equation}
\begin{aligned}
    -\frac{\partial p_\infty(z,t)}{\partial x} 
    &= \frac{d{G_m}_x(z,t)}{dt} 
    - f_c\,{G_m}_y(z,t) 
    - \frac{\partial \langle \mathcal{T}_{xz}(z,t) \rangle}{\partial z}, \\
    -\frac{\partial p_\infty(z,t)}{\partial y} 
    &= \frac{d{G_m}_y(z,t)}{dt}
    + f_c\,{G_m}_x(z,t) 
    - \frac{\partial \langle \mathcal{T}_{yz}(z,t) \rangle}{\partial z},
\end{aligned}
\label{eq:geostrophicBalanceUnsteady3}
\end{equation}
%
where the last term in the right-hand side represents the frictional drag force, incorporating both the resolved-scale Reynolds stress and the subgrid-scale contribution modeled in LES as
\begin{equation}
\mathcal{T}_{i3}(z) = -R_{i3}(z)  - \mathcal{D}_{i3}(z) - \tau_{i3}^{\mathrm{sgs}}(z),
\label{eq:totalShearStress}
\end{equation}
where \(i = 1,2\) corresponds to the horizontal directions \(x\) and \(y\), respectively. Here, \(R_{i3}(z) = \langle \overline{\bar{u}_i'\,\bar{u}_3'} \rangle\) denotes the resolved Reynolds shear stress, $\mathcal{D}_{i3}(z) = \bigl\langle \bar{u}_i''\,\bar{u}_3''\bigr\rangle$ is the dispersive stress component and \(\tau_{i3}^{\mathrm{sgs}}(z)\) represents the subgrid-scale shear stress modeled in LES. Since the present study considers homogeneous flat terrain, dispersive stresses are negligible and are therefore omitted in the computation of \(\mathcal{T}_{i3}(z)\). 

Horizontal temperature gradients, which are responsible for vertically varying geostrophic winds, act on spatial scales much larger than the domain sizes considered in this study. For instance, in the mid-latitude troposphere, a typical horizontal temperature gradient of approximately \(1\,\mathrm{K}\) per \(100\,\mathrm{km}\) can induce a geostrophic wind shear of nearly \(3.5\,\mathrm{m\,s^{-1}}\) per \(\mathrm{km}\)~\citep{sorbjan2004large}. Consequently, horizontal temperature gradients are not modelled in this study and the effect of baroclinicity is directly represented by using mesoscale wind components \( G_{x}(z,t) \) and \( G_{y}(z,t) \), typically obtained from numerical weather prediction (NWP) models or field observations. Using the subscript \( m \) to denote mesoscale data, the inertial and Coriolis terms in Eq.~\ref{eq:geostrophicBalanceUnsteady3} are derived from these mesoscale wind fields. In contrast, the frictional drag force term is computed directly from the microscale LES solver using horizontal plane averaging and temporal filtering, where the filtering window is chosen to match the temporal resolution of the mesoscale data. 

HGWPA is a hybrid methodology in which the WPA approach is used to compute the virtual potential temperature source term \( F_\theta \), while the momentum source term \( F_{u_i} \) is approximated using the large-scale pressure gradient computed from Eq.~\ref{eq:geostrophicBalanceUnsteady3}. Although WPA is used in this study, previous methods such as IPA or DPA can also be employed for computing the temperature source term. However, we found that WPA integrates particularly well with the geostrophic profile assimilation compared to IPA or DPA, yielding improved accuracy in predicting the mean and turbulent fields. HGWPA enables a more physically consistent specification of the large-scale pressure gradient in LES of the ABL by accounting for baroclinic variability, temporal evolution, and turbulence-induced momentum transfer. In contrast, other methods such as DPA, IPA, and WPA introduce forcing in Eq.~\ref{eq:momentumCartesian} based on the error between mesoscale data and microscale solver outputs. Consequently, errors stemming from the mismatch between mesoscale forcing and LES data may affect the mean and turbulent flow fields. Although IPA and WPA introduce vertical coherence in the error term using a curve-fitting procedure, this fitting is applied independently at each time step. As a result, these approaches lack temporal coherence in the applied forcing. This limitation motivates the use of a physics-based approach such as HGWPA, which overcomes these deficiencies by directly incorporating the relevant physical processes into the computation of the large-scale pressure gradient.

To evaluate the relative importance of the frictional drag term within the ABL when modeling the large-scale pressure gradient, three different cases are considered here as illustrated in Fig.~\ref{fig:HGWPACaseSetup}. In Case 1 termed HGWPA1, the frictional drag term is neglected in the computation of the large-scale pressure gradient. 
A vertically uniform pressure gradient 
is assumed within the ABL with its value taken as that just above the ABL. Above the ABL, the large-scale pressure gradient is computed using Eq.~\ref{eq:geostrophicBalanceUnsteady2}. Case 2 (HGWPA2) includes the 
frictional drag term and the large-scale 
pressure gradient is computed throughout the domain using 
Eq.~\ref{eq:geostrophicBalanceUnsteady3}. Case 3 (HGWPA3) is similar to Case 2 such that the frictional drag term is included using 
Eq.~\ref{eq:geostrophicBalanceUnsteady3}. However, in this case, the effect of the frictional drag is restricted to within the ABL by multiplying it with a scaling function $f_s(z)$ defined as 
\begin{equation}
	f_s(z) = \frac{1}{2}\left[1 - \tanh\left(\frac{7(z-H)}{\Delta}\right)\right], 
	\label{eq:geoDampingFunction}
\end{equation}
where \(H\) is the ABL height and corresponds to the level at which the scaling reduces to half its maximum value. The parameter \(\Delta\), set to \(0.2H\), defines the transition width over which the frictional drag term smoothly decays to 0 above the ABL. 
\begin{figure}
    \centering
    {\large (a) HGWPA1} \\[2pt]
    \includegraphics[width=0.6\linewidth]{Plots/geostrophicsetup1.pdf} \\[2pt] 
    \begin{tabular}{cc}
    {\large (b) HGWPA2} & {\large (c) HGWPA3} \\ [2pt]
         \includegraphics[width=0.46\linewidth]{Plots/geostrophicsetup2.pdf} &
         \includegraphics[width=0.46\linewidth]{Plots/geostrophicsetup3.pdf}
    \end{tabular}    
\caption{Schematic illustration of the momentum forcing techniques for hybrid geostrophic-wavelet approaches. The instantaneous ABL height, shown by the purple dashed line, is used in HGWPA1 and HGWPA3. In HGWPA1, constant geostrophic forcing is applied below the ABL height, while variable geostrophic forcing derived from mesoscale data is applied above, including only the inertial term in geostrophic balance (Eq.~\ref{eq:geostrophicBalanceUnsteady2}). HGWPA2 (Eq. \ref{eq:geostrophicBalanceUnsteady3}) incorporates turbulent drag forcing without the need to compute the ABL height; the forcing naturally adapts to the ABL structure. HGWPA3 (Eq.~\ref{eq:geostrophicBalanceUnsteady3}) is similar to HGWPA2 but applies a tapering of the drag force above the ABL to zero using a hyperbolic tangent function (red curve), as defined in Eq.~\ref{eq:geoDampingFunction}.}
\label{fig:HGWPACaseSetup}
\end{figure}

Cases 1 and 3 require the ABL height to be computed dynamically over the diurnal cycle of the simulation. The instantaneous ABL height \(H\), is estimated based 
on two criteria: the turbulent shear stress and the turbulent heat flux~\citep{allaerts2015large}. Under stable conditions, \( H \) is defined as the height at which the total turbulent 
shear stress defined using Eq.~\ref{eq:totalShearStress} decreases below \(5\%\) of the surface stress. Under convective conditions, 
\( H \) is taken as the height where the turbulent heat flux, defined as:
\begin{equation}
q_{z}(z) 
= \bigl\langle \bar{w}'\,\bar{\theta}' \bigr\rangle(z)
\; + \;\bigl\langle q_{z}^{\mathrm{sgs}}(z)\bigr\rangle,
\label{eq:geoABL}
\end{equation}
reaches its minimum, where \( q_{z}^{\mathrm{sgs}}(z) \) denotes the subgrid-scale contribution to the vertical 
heat flux. A transition between the shear-based and 
heat-flux-based definitions is applied during the dirunal cycle depending on the sign of the time-filtered surface heat flux; a heat-flux-based ABL height definition is used when this quantity is positive, and a shear-stress-based ABL height definition is used when it is negative.

%%%%%%%%%%%%%%%%%%%%%%%%%%%%%%%%%%%%%%%%%%%%%%%%%%%%%%%%%%%%%
%%%%%%%%%%%%%%%%%%%% Model Setup %%%%%%%%%%%%%%%%%%%%%%%%%%%%
%%%%%%%%%%%%%%%%%%%%%%%%%%%%%%%%%%%%%%%%%%%%%%%%%%%%%%%%%%%%%
\section{Observational data and model setup}
\label{sec:computationalModel}
%
To demonstrate our MMC approach, we utilize data from the American WAKE experimeNt (AWAKEN) field campaign in northern Oklahoma, USA. The AWAKEN project investigates the impact of large-scale wind farm wakes and their interaction with the atmospheric boundary layer in a region that includes multiple wind farms. The flow field in this region has been characterized through an extensive field campaign with several instrumented sites \citep{moriarty2023awaken}. Among these sites, A2 (see Fig. \ref{fig:A2Bsite}) is selected as the primary instrumentation location for southerly winds upstream of the King plains wind farm and is used to extract the time-height velocity flow field for this study. Site A2 is instrumented with a scanning lidar, surface sonic anemometer, and surface meteorological station. The scanning lidar provides $10$-minute averaged measurements of the velocity profile and second-order wind statistics from $100~\mathrm{m}$ to $3000~\mathrm{m}$ above ground. However, during the period utilized in this  study, data was only available up to $1100~\mathrm{m}$ above ground. Momentum and heat flux measurements are additionally obtained from the surface sonic anemometer mounted at 4 m above ground at site A2, to compute the surface Obukhov length ($L$). For validation, horizontal wind speed and wind direction data from virtual tower measurements at site A1 are also utilized. These data are available at 3-min intervals and span heights from 15 m to 325 m, derived through dual-Doppler analysis of two scanning lidars located at sites A5 and A7.
%
\begin{figure}[htb!]
    \centering
    \includegraphics[width=0.9\linewidth]{Plots/observationalCP.pdf}
\caption{Layout of the AWAKEN instrumented sites.
The left panel shows all instrumented sites of the AWAKEN project, while the right panel focuses on the area of interest for the present MMC study. All mesoscale data were obtained from the A2 site, which hosts a scanning lidar instrument. Black dots enclosed in red circles represent turbines instrumented as part of the AWAKEN campaign. The white star represents the remote sensing and surface instruments, whereas the blue and purple stars denote the X-band radar and the meteorological tower, respectively. The positions of all turbines are not shown for clarity, but are available in \citet{moriarty2023awaken}.}


\label{fig:A2Bsite}
\end{figure}

Virtual potential temperature vertical profile and surface temperature data are obtained from $30$-minute averaged measurements using the NREL ASSIST-II (SN 11) infrared spectrometer located at site G (see Fig. \ref{fig:A2Bsite}). The instrument provides continuous thermodynamic data from the surface up to an altitude of $17~\mathrm{km}$. 

The simulation period spans the diurnal cycle from September $1$ to September $2$, $2023$. This period for the case study is chosen based on the combined availability of the lidar and thermodynamic profiler data at both the inflow sites. Additional constraints are imposed by the requirement of a southerly flow (between $150^\circ$ and $210^\circ$) to ensure that the flow at site A2 is not waked by other wind farms. Furthermore, the data period was carefully screened to identify intervals with relatively consistent wind speed and direction, alongside the development of a stable boundary layer characterized by the formation of a nocturnal low-level jet.

\begin{figure}[htb!]
    \centering
    \begin{tabular}{cc}
         {\large (a) {Observational-U}} & {\large (d) {Observational-V}}  \\
         \includegraphics[width=0.45\linewidth]{Plots/ObservedU.png} &
         \includegraphics[width=0.45\linewidth]{Plots/ObservedV.png} \\
        {\large (b) {WRF-U}} & {\large (e) {WRF-V}} \\
         \includegraphics[width=0.45\linewidth]{Plots/WRFU.png} &
         \includegraphics[width=0.45\linewidth]{Plots/WRFV.png} \\
        {\large (c) {Blended-U}} & {\large (f) {Blended-V}} \\
         \includegraphics[width=0.45\linewidth]{Plots/BlendedU.png} &
         \includegraphics[width=0.45\linewidth]{Plots/BlendedV.png} \\
    \end{tabular}
\caption{Time-height contour of horizontal wind speed components U and V from observations, WRF simulations, and blended data. Observational data are available only between $100$ and $1100~\mathrm{m}$ (panels a, d). To extend the profiles above this range, mesoscale fields from WRF (panels b, e) are used. The blended datasets (panels c, f) combine the observational data within the available height range with WRF outputs above.}

    \label{fig:dpacases}
\end{figure}

We conduct MMC simulations for the full diurnal cycle to evaluate the performance of the direct and indirect profile assimilation algorithms and to test our newly proposed mesoscale forcing methods. All simulations are performed using the LES code TOSCA (Toolbox fOr Stratified Convective Atmospheres), which solves the filtered incompressible Navier–Stokes equations with Coriolis force and the Boussinesq approximation for the buoyancy term, using a finite-volume method expressed in a generalized curvilinear coordinates \citep{stipa2023tosca}. Additionally, thermal stability effects in the ABL and free atmosphere are accounted by solving a potential temperature transport equation. TOSCA is optimized for large-scale simulations and makes use of parallel libraries such as PETSc \citep{balay2022petsc} and HYPRE \citep{falgout2002hypre}, along with high-performance I/O through the HDF5 library \citep{HDFGroup2025}. The computational domain is horizontally periodic and spans $3~\mathrm{km} \times 3~\mathrm{km} \times 2~\mathrm{km}$, with a uniform grid spacing of $\Delta x = \Delta y = 20~\mathrm{m}$ in the horizontal directions and $\Delta z = 10~\mathrm{m}$ in the vertical direction, giving a consistent cell aspect ratio of $2$. This grid resolution is chosen based on the grid study described in Appendix~\ref{sec:GRAstudy}.

The LES is initialized at $0600~\mathrm{UTC}$ ($0000~\mathrm{LT}$) on September $1$, $2023$, and initially run with steady momentum and temperature forcing for $4~\mathrm{hr}$ to allow the turbulence to spin up, before profile assimilation begins for a diurnal cycle. Then, the simulation is advanced for $30~\mathrm{hr}$ with a time step of $1~\mathrm{s}$. Standard wall stress and heat flux boundary conditions \citep{schumann1975subgrid, grotzbach1987direct} are applied at the surface, with a uniform aerodynamic roughness length of $0.15~\mathrm{m}$--an effective aerodynamic roughness value for the region where the observational data were collected \citep{cheung2023investigations} and surface temperature obtained from the NREL ASSIST-II (SN 11) infrared spectrometer located at site G.

We also conducted a WRF simulation to compare the proposed profile assimilation forcing terms with the advection and thermal tendencies derived from the mesoscale solver. The mesoscale atmospheric flow is simulated using the WRF version 4.1.3 \citep{skamarock2004evaluating}. WRF solves the fully compressible, non-hydrostatic Navier-Stokes equations and uses a complete suite of physical parameterization schemes for subgrid scale processes. The WRF model uses the ERA5 reanalysis data sets \citep{hersbach2020era5} for initial and boundary conditions. ERA5 provides hourly, global atmospheric, land-surface and sea-state gridded data on a $0.25^\circ$ resolution. WRF is dynamically downscaled with three domains with a grid spacing of $9$, $3$ and $1$ $\mathrm{km}$, respectively. It is set with a total of $40$ vertical levels and a vertical spacing of $20~\mathrm{m}$ near the surface, including approximately $10$ vertical levels in the lowest $200~\mathrm{m}$. The simulation is performed with a $24$-hour spin-up and spectral nudging is turned on to selectively nudge the largest wavelengths of e.g. U, V and geopotential height to the ERA5 data above the PBL \citep{miguez2004spectral}. Model outputs are saved every 10 minutes. 

Figure~\ref{fig:dpacases} shows the comparison between observational data and WRF simulations for the horizontal wind components. Although the WRF results are smoothed compared to observations, they capture the overall diurnal trends. Since the observational data extends only up to $1100~\mathrm{m}$, they are blended with the WRF outputs to produce the combined dataset shown in Fig.\ref{fig:dpacases}c and \ref{fig:dpacases}f. This blended dataset is used for the hybrid geostrophic-wavelet simulations (Sect.\ref{sec:hybridPA}), while the observational dataset in Fig.\ref{fig:dpacases}a and \ref{fig:dpacases}d is used for the error-based profile assimilation simulations (Sect.~\ref{sec:errorBased}), with constant forcing applied outside the observational range (below 100 m and above 1100 m). Using the blended dataset in the hybrid geostrophic-wavelet cases ensures data availability above the boundary layer height to enforce geostrophic balance, particularly in the convective regime where the boundary layer height can exceed 1100 m.
%%%%%%%%%%%%%%%%%%%%%%%%%%%%%%%%%%%%%%%%%%%%%%%%%%%%
%%%                   Results                    %%%
%%%%%%%%%%%%%%%%%%%%%%%%%%%%%%%%%%%%%%%%%%%%%%%%%%%%
\section{Results}
\label{sec:results}

To facilitate a comprehensive analysis of the proposed assimilation approaches, we organize the analysis into two main categories. The first category comprises error-based methods: DPA, IPA, and WPA. The second includes hybrid geostrophic-wavelet approaches: HGWPA1, HGWPA2, and HGWPA3. In Sect.~\ref{sec:forcingTerms}, we compare the forcing terms obtained from the various methods. Next, Sects.~\ref{sec:errorBased} and \ref{sec:hybridPA} present analyses of the diurnal evolution of mean and turbulence statistics and the vertical structure of the boundary layer for the error-based and hybrid profile assimilation methods, respectively.

\subsection{\textbf{Comparison of forcing terms}}
\label{sec:forcingTerms}

Profile assimilation methods evaluate the time- and height-varying momentum forcing term \(F_{u_i}\) and the virtual potential temperature forcing term \(F_{\theta}\) in the microscale solver (Eqs. \ref{eq:momentumCartesian} and \ref{eq:temperatureCartesian}). This section analyzes how these source terms vary when computed using error-based assimilation methods (DPA, IPA, WPA) and geostrophic-wavelet hybrid approaches (HGWPA1, HGWPA2, HGWPA3). The momentum and temperature forcing terms, denoted by \(F_{u_i}\) and \(F_{\theta}\), are calculated using Eqs. \ref{dpaForcing}, \ref{ipaForcing}, and \ref{wpaForcing} for DPA, IPA, and WPA, respectively. For HGWPA1, the forcing is derived from Eq. \ref{eq:geostrophicBalanceUnsteady2}, while HGWPA2 and HGWPA3 use Eq. \ref{eq:geostrophicBalanceUnsteady3}. For reference, forcing terms from the mesoscale model (WRF) are used for comparison.

Figures~\ref{fig:momSrcVert}(a--c) compare the vertical profiles of the $x$-momentum component, $y$-momentum component and virtual potential temperature forcing terms respectively obtained from the various profile assimilation methods. The height- and time-varying forcing data were extracted from the microscale solver and averaged over one-hour windows centered at 1730--1830 UTC, 2330--0030 UTC, and 0530--0630 UTC on the second day. These intervals correspond to convective (\(L = -23.97\)~m), transitional (\(L = -330.19\)~m), and stable (\(L = 15.78\)~m) boundary-layer regimes respectively, where \(L\) is the surface Obukhov length computed from sonic-anemometer heat-flux and momentum flux measurements.

\begin{figure}[!htbp]
\centering
    \begin{tabular}{c}
        {\large (a)} \\ [2pt]
        \includegraphics[width=0.75\textwidth]{Plots/srcTerms/momSrcX_profiles_new.png}  \\ [2pt]
        {\large (b)} \\[2pt]
        \includegraphics[width=0.75\textwidth]{Plots/srcTerms/momSrcY_profiles_new.png} \\ [2pt]
        {\large (c)} \\
          \includegraphics[width=0.75\textwidth]{Plots/srcTerms/momSrcT_profiles_new.png} \\ [2pt]
    \end{tabular}    
\caption{Vertical profiles of mesoscale forcing (a) $x$-momentum component, (b) $y$-momentum component, (c) virtual potential temperature, averaged over a one-hour time period. Profiles are shown with averaging about 1800 UTC (left panels), 0000 UTC (middle panels), and 0600 UTC (right panels) for the first and second September 2023. Forcing terms are computed using DPA, IPA, WPA, and three cases of HGWPA. Additionally, the momentum terms have been scaled with the inverse of the Coriolis force $f_c$.}
\label{fig:momSrcVert}
\end{figure}

Across all stability regimes (Fig. \ref{fig:momSrcVert}), the DPA method exhibits pronounced fluctuations in both momentum and temperature forcing terms within the boundary layer, indicating limited vertical coherence. However, as turbulence weakens during the stable regime (Fig. \ref{fig:momSrcVert} right panels), DPA forcing gradually converges toward the profiles obtained from the other assimilation methods, above the boundary layer. Similar trends have been reported by \citet{allaerts2020development}.

In contrast, IPA and WPA yield smoother and more vertically coherent profiles in the momentum forcing components ($F_u$ and $F_{v}$), as shown in Figs. \ref{fig:momSrcVert}a and \ref{fig:momSrcVert}b. Both of these methods are effective in suppressing the noise previously observed with DPA. However, a clear difference is observed in the case of virtual potential temperature, as shown in Fig.~\ref{fig:momSrcVert}c, where the wavelet-based algorithm captures local variations and maintains vertical correlation. In contrast, IPA leads to over-smoothing and is limited by the cubic polynomial fitting, which limits local variations and fails to capture the complex vertical structure. This finding underscores a key limitation of IPA: its ability to represent complex vertical structure is strongly dependent on the chosen polynomial order with higher-order polynomials intensifying the boundary edge runaway effect, i.e., larger source terms near the boundary \citep{allaerts2020development}. Therefore, we choose a wavelet-based profile assimilation method to combine with geostrophic balance-based forcing methods in the hybrid strategy (Sect.~\ref{sec:hybridPA}).

Figures~\ref{fig:momSrcVert}(a--c) also illustrate the performance of the hybrid forcing schemes (HGWPA1-3). In all three hybrid methods, the virtual potential temperature forcing is treated using the wavelet-based method, although slight differences in the temperature forcing vertical profiles arise due to variations in the momentum forcing formulations. HGWPA1 includes only the unsteady geostrophic term and applies a constant forcing within the ABL, enforcing geostrophic balance above it. HGWPA2 and HGWPA3 incorporate an additional turbulent drag term, with the key difference being that HGWPA3 estimates the ABL height and gradually tapers the drag force to zero above it, while HGWPA2 applies the drag throughout without explicitly defining the ABL top.
These differences are most evident under stable conditions (at 06:00 UTC), where HGWPA2 exhibits greater variability with height above the ABL due to the development of a residual layer---a remnant of the daytime convective boundary layer---characterized by enhanced shear and intermittent turbulence. HGWPA3, in contrast, produces smoother vertical transitions near the ABL top, aligning closely with HGWPA1 above the boundary layer where only the unsteady geostrophic term remains active. Near the surface, both HGWPA2 and HGWPA3 show sharper vertical gradients in the forcing terms, attributable to the inclusion of the turbulent drag component.

It is notable that, despite being based on fundamentally different approaches---error-fitting versus physically derived schemes---the two categories generate source terms of similar magnitude across most of the boundary layer under the different ABL regimes.

\begin{figure}[htb!]
    \centering
    \begin{tabular}{cc}
        {\large (a)} & {\large (b)}  \\[5pt]
        \includegraphics[width=0.48\textwidth]{Plots/srcTerms/timeHeightTemp_DPA_IPA_WPA.png} &
      \includegraphics[width=0.48\textwidth]{Plots/srcTerms/timeHeightTemp_HGWPA.png} \\
    \end{tabular}
    \caption{Time--height profiles of mesoscale virtual potential temperature forcing term computed from data assimilation between 1-2 September 2023. Forcing terms are computed using DPA, IPA, WPA and the three cases of HGWPA. Black line represents the mean ABL height filtered using a 1 hour moving time window.}
    \label{fig:momSrcSurfaceWPA}
\end{figure}

For a qualitative comparison, Fig.~\ref{fig:momSrcSurfaceWPA}(a--b) shows the time-height contour plots of the virtual potential temperature forcing terms, overlaid with the ABL height computed during the diurnal evolution of the boundary layer. All figures indicate a gradual growth of the boundary layer height during the transition from stable to convective conditions, followed by an abrupt drop in ABL height when shifting from convective to stable regimes. IPA predicts a higher boundary layer height than the other methods, with a steeper gradient during the transition from stable to convective conditions. Regarding the virtual potential temperature forcing term, DPA shows a noisy forcing signal, particularly during the daytime convective period. WPA smooths the forcing term, correctly capturing the strong heating event during the convective boundary layer. The HGWPA methods show similar trend to WPA as it uses the wavelet method for the virtual potential temperature forcing. In contrast, IPA tends to oversmooth temporal and vertical variations, masking transient features, although it captures the general trends in both stable and convective regimes.

\FloatBarrier 
\subsection{\textbf{Error-based profile assimilation}}
\label{sec:errorBased}

\subsubsection{Diurnal evolution of data driven LES}
This section presents a comparative analysis of three data assimilation techniques: direct, indirect, and wavelet-based. The evaluation focuses on how effectively each method captures the diurnal evolution of key flow variables. Mean-flow quantities, including horizontal wind speed ($S=\sqrt{U^2 + V^2}~\mathrm{m/s}$), wind direction ($D$), and virtual potential temperature ($\theta_v$), are considered alongside turbulent quantities such as shear stress ($\tau_z$) and turbulent kinetic energy ($k$). Here, we evaluate the shear stress \(\tau_z\) and the turbulent kinetic energy from the resolved-scale Reynolds stress tensor \(R_{ij} = \overline{u_i' u_j'}\). The TKE is defined as \(k = \frac{1}{2} \operatorname{Tr}(R_{ij}) = \frac{1}{2}(R_{11} + R_{22} + R_{33})\), where \(\operatorname{Tr}(R_{ij})\) denotes the trace of the Reynolds stress tensor. The vertical shear stress is computed as \(\tau_z = \sqrt{R_{13}^2 + R_{23}^2}\), representing the magnitude of the vertical momentum transport due to streamwise and spanwise velocity fluctuations. These variables are assessed using time series and time-height contour plots, with comparisons made against observational data from scanning lidar and virtual tower data to evaluate how realistically each method reproduces the physical fields of interest. The observational data are averaged over $10$-$\mathrm{min}$, while the microscale flow quantities represent planar averages and are also temporally averaged in $10$-$\mathrm{min}$ bins.

Since the IPA is susceptible to boundary edge effects near the ground, we compare results at two different heights:  ($100~\mathrm{m}$, see Figs.~\ref{fig:timehistoryMeanDPA100m} and \ref{fig:timehistoryTurbDPA100m}) and at $200~\mathrm{m}$ (see Figs.~\ref{fig:timehistoryMeanDPA200m} and \ref{fig:timehistoryTurbDPA200m}). These heights correspond to the hub height of modern wind turbines and lie within the surface layer, where turbulent phenomena dominate and both mean and turbulent quantities are critical for understanding the flow physics in realistic scenarios.

\begin{figure}[htb!]
    \centering
    \begin{tabular}{c}
        {\large{(a)}} \\ [2pt]
         \includegraphics[width=0.6\linewidth]{Plots/fullDomain/windSpeed_DPA_100m.png} \\
        {\large{(b)}} \\
        \includegraphics[width=0.6\linewidth]{Plots/fullDomain/windDirection_DPA_100m.png} \\
    \end{tabular}
    \caption{Temporal variation of (a) horizontal wind speed at $100~\mathrm{m}$ height and its zoomed-in section (0900--1500 UTC),  (b) wind direction, and virtual potential temperature at $100~\mathrm{m}$.}
    \label{fig:timehistoryMeanDPA100m}
\end{figure}

\begin{figure}[htb!]
    \centering
    \begin{tabular}{c}
        {\large{(a)}}  \\ [2pt]
         \includegraphics[width=0.6\linewidth]{Plots/fullDomain/kineticEnergy_DPA_100m.png} \\
        {\large{(b)}} \\
        \includegraphics[width=0.6\linewidth]{Plots/fullDomain/shearStress_DPA_100m.png} \\
    \end{tabular}
    \caption{Temporal variation of (a) turbulent kinetic energy and its zoomed-in section (1500--2100 UTC),  (b) shear stress, and its detailed section (1500--2100 UTC) at $100~\mathrm{m}$ height.}
    \label{fig:timehistoryTurbDPA100m}
\end{figure}


\begin{figure}[htb!]
    \centering
    \begin{tabular}{c}
        {\large{(a)}}  \\ [2pt]
         \includegraphics[width=0.6\linewidth]{Plots/fullDomain/windSpeed_DPA_200m.png} \\
        {\large{(b)}} \\
        \includegraphics[width=0.6\linewidth]{Plots/fullDomain/windDirection_DPA_200m.png} \\
    \end{tabular}
    \caption{Temporal variation of (a) horizontal wind speed and its zoomed-in section (0900--1500 UTC), (b) wind direction, and virtual potential temperature, at $200~\mathrm{m}$ height.}
    \label{fig:timehistoryMeanDPA200m}
\end{figure}

\begin{figure}[htb!]
    \centering
    \begin{tabular}{c}
        {\large{(a)}}  \\ [2pt]
         \includegraphics[width=0.6\linewidth]{Plots/fullDomain/kineticEnergy_DPA_200m.png} \\
        {\large{(b)}} \\
        \includegraphics[width=0.6\linewidth]{Plots/fullDomain/shearStress_DPA_200m.png} \\
    \end{tabular}
    \caption{Temporal variation of (a) turbulent kinetic energy and its zoomed-in section (1500--2100 UTC),  (b) shear stress, and its detailed section (1500--2100 UTC), at $200~\mathrm{m}$ height.}
    \label{fig:timehistoryTurbDPA200m}
\end{figure}

Figures~\ref{fig:timehistoryMeanDPA100m}a and \ref{fig:timehistoryMeanDPA100m}b show the diurnal evolution of horizontal wind speed, wind direction, and virtual potential temperature estimated by the three assimilation methods at $100~\mathrm{m}$. Both the direct and wavelet-based data assimilation methods capture the wind speed trends and remain consistent with the observational data. In contrast, the IPA method shows a noticeable deviation with the measurements during the stable period between 0900 and 1500 UTC. A zoomed-in view of the wind speed in Fig.~\ref{fig:timehistoryMeanDPA100m}a highlights the difference between IPA and the measured values. This behavior stems from the cubic polynomial used in the IPA formulation, which induces over-relaxation of the error term compared to WPA, leading to a larger deviation from the mesoscale data. It may also reflect an influence of the boundary edge runaway effect. For wind direction and virtual potential temperature, all assimilation methods generally follow the observational data throughout the diurnal cycle. 

Figures~\ref{fig:timehistoryTurbDPA100m}a and \ref{fig:timehistoryTurbDPA100m}b illustrate the diurnal evolution of turbulent flow quantities at $h=100~\mathrm{m}$. Unlike velocity and virtual potential temperature, turbulence quantities are not directly constrained by error forcing but instead emerge as a consequence of the energy cascade from the mesoscale to the LES scale. It is therefore important to assess how well the different methods reproduce turbulence characteristics compared to the observational data. Figure~\ref{fig:timehistoryTurbDPA100m}a shows the time history of turbulent kinetic energy, which indicates that the direct profile assimilation algorithm artificially amplifies turbulent kinetic energy, overpredicting TKE during the transition and convective periods. This overprediction leads to excessive energy accumulation, which appears to disrupt the natural energy cascade, as insufficient dissipation fails to balance the energy production. On the other hand, the IPA algorithm underestimates the TKE during the convective regime compared to the observational data. In contrast, the wavelet-based profile assimilation algorithm estimates the TKE across the entire diurnal cycle reasonably well and generally follows the measured TKE data. Fig.~\ref{fig:timehistoryTurbDPA100m}a also provides a magnified view of TKE, highlighting the clear distinction among the methods.

Figure~\ref{fig:timehistoryTurbDPA100m}b shows the time history of shear stress at \( h = 100~\mathrm{m} \), where the indirect, direct, and wavelet-based data assimilation methods are used to estimate the shear stress. The results show distinct differences in how each method captures coherent structures in the convective boundary layer. A ramp-like pattern in the shear stress time series indicates the presence of coherent turbulent structures or organized events in the atmospheric boundary layer. These features are common in the convective boundary layer, especially during the daytime, when surface heating causes buoyancy-driven turbulence in the form of thermal plumes~\citep{gao1989observation,Paw1992,vidakovic2009statistical,li2011coherent}. The IPA method over-smooths the shear stress profile and fails to resolve the ramp-like structures. It struggles to distinguish between organized and unorganized events due to the local and non-periodic nature of coherent structures in the mesoscale data. These structures contribute significantly to spectral and co-spectral energy at scales smaller than their characteristic size, making them difficult to isolate using classical filtering approaches such as box filtering, Gaussian filtering, or polynomial fitting. As noted by~\citet{allaerts2020development}, Gaussian filters can even produce incorrect statistics in such cases.

The DPA method does not apply any filtering, which allows both coherent and incoherent modes to remain in the signal. This leads to overprediction of the shear stress magnitude. In contrast, the wavelet-based method adapts to the localized, non-periodic nature of coherent structures. It uses the MRA technique to remove uncorrelated stochastic noise while preserving dominant features of the shear stress, resulting in a shear stress profile that closely matches observational data. Fig.~\ref{fig:timehistoryTurbDPA100m}b also provides a zoomed-in view of the time history, showing that the wavelet method better resolves these localized, non-periodic features.

At $h=200~\mathrm{m}$, Figs.~\ref{fig:timehistoryMeanDPA200m}a and \ref{fig:timehistoryMeanDPA200m}b reveal that the mean profiles, particularly wind speed, exhibit improved convergence towards observational data with the IPA algorithm, compared to $h=100~\mathrm{m}$. The diurnal evolution of other mean profiles, such as wind direction and virtual potential temperature, shows no significant differences and follows similar trends to the $h=100~\mathrm{m}$ profiles. Moreover, turbulent flow quantities like turbulent kinetic energy and shear stress at $h=200~\mathrm{m}$, (Figs.~\ref{fig:timehistoryTurbDPA200m}a and \ref{fig:timehistoryTurbDPA200m}b) also show no improvement with either DPA or IPA, both struggling to capture the intermittent and organized events at this height. The wavelet-based data assimilation method is again able to reproduce the mean and turbulent flow quantities more accurately.

The wavelet-based data assimilation method consistently outperforms other algorithms in capturing both mean and turbulent flow quantities at $h=100~\mathrm{m}$ and $h=200~\mathrm{m}$. This superior performance stems from the wavelet method's ability to dynamically adapt to localized flow structures, enabling it to capture the dominant and localized turbulent features and coherent structure signatures that other methods, such as IPA and DPA, fail to resolve. 

In Fig.~\ref{fig:timeheightcontourDPA}, time-height contours of horizontal wind speed, wind direction, virtual potential temperature, TKE, shear stress, and the virtual potential temperature gradient are presented for DPA, IPA, and WPA, alongside observations from the scanning lidar. The estimated ABL height is overlaid to facilitate analysis of the ABL vertical structure throughout the diurnal cycle.

Figures~\ref{fig:timeheightcontourDPA}a and \ref{fig:timeheightcontourDPA}b show that all methods successfully reproduce the overall evolution of wind speed and wind direction. As the ABL transitions from stable to convective and back to stable conditions, a low-level jet develops during the stable period, disappears under convective conditions due to enhanced mixing, and re-emerges during the subsequent stable phase as wind speeds increase. Closer inspection reveals differences in how the three methods represent the low-level jet. DPA predicts wind speeds most accurately, as expected from the direct forcing of the error term. Between IPA and WPA, IPA produces a broader and more diffuse jet with lower peak velocities, whereas WPA yields a narrower jet core with higher velocities, more consistent with the observational data. 

Figures~\ref{fig:timeheightcontourDPA}d and \ref{fig:timeheightcontourDPA}e show the time–height contours of TKE and turbulent shear stress. During the diurnal cycle, elevated TKE and shear stress are observed in the convective regime, whereas both remain much lower in the stable regime, with intermittent bursts of higher values. DPA produces substantially larger TKE during the convective period (clipped at 3 $m^2/s^2$) and elevated values even under stable conditions. In contrast, IPA underpredicts both TKE and shear stress, while WPA results agree most closely with the observational data. Overall, WPA provides a good compromise, accurately representing wind speed and wind direction while also capturing the turbulent characteristics most effectively.

Figure~\ref{fig:timeheightcontourDPA}f shows the vertical virtual potential temperature gradient over the diurnal cycle. During the daytime convective regime, the growth of the boundary layer is evident through the upward rise of the entrainment zone, marked by a positive temperature gradient. Within the mixed layer, the gradient remains near zero, while a region of negative virtual potential temperature gradient develops close to the surface. This feature is reproduced more accurately by the WPA method than by the IPA method and aligns more closely with observations. Following sunset, the transition to a stable boundary layer becomes evident, marked by the emergence of a positive near-surface temperature gradient between 0000 and 0600 UTC.

\begin{figure}[htb!]
    \centering
    \begin{tabular}{cc}
        \large{(a)} & \large{(d)} \\
        \includegraphics[width=0.45\linewidth]{Plots/fullDomain/timeHeight_windspeed.png} &
        \includegraphics[width=0.45\linewidth]{Plots/fullDomain/timeHeight_tke.png} \\
         \large{(b)} & \large{(e)} \\
        \includegraphics[width=0.45\linewidth]{Plots/fullDomain/timeHeight_winddirection.png} &
        \includegraphics[width=0.45\linewidth]{Plots/fullDomain/timeHeight_shearstress.png} \\
         \large{(c)} & \large{(f)} \\
        \includegraphics[width=0.48\linewidth]{Plots/fullDomain/timeHeight_theta.png} &
        \includegraphics[width=0.45\linewidth]{Plots/fullDomain/timeHeight_dtdz.png}
    \end{tabular}
\caption{Time-height contour plots derived from DPA, IPA, and WPA profile assimilation methods, compared with mesoscale observational data. Panels (a) horizontal wind speed, (b) wind direction, (c) virtual potential temperature, (d) turbulent kinetic energy, (e) shear stress, and (f) virtual potential temperature gradient. Each plot is overlaid with boundary layer heights computed via the corresponding assimilation techniques.}
    \label{fig:timeheightcontourDPA}
\end{figure}

\subsubsection{Vertical structure of the atmospheric boundary layer}

This section examines the evolution of the vertical structure of the atmospheric boundary layer over a diurnal cycle. We present vertical profiles of mean and turbulent flow quantities estimated using error-based profile assimilation methods. The microscale vertical data are one-hour averages centered at 1800~UTC (convective boundary layer), 0000~UTC (transition phase), and 0600~UTC (stable boundary layer) on the second day with further spatial averaging over horizontal planes. Additionally, the observational data are also time-averaged over a one hour period for comparison.  

Figures~\ref{fig:verticalFullDomain1}a and \ref{fig:verticalFullDomain1}b show the vertical profiles of wind speed, wind direction, virtual potential temperature, and turbulent kinetic energy for the convective and transitional regimes, respectively. At 1800~UTC (Fig.~\ref{fig:verticalFullDomain1}a), IPA slightly overestimates wind speed within the surface layer (ABL height~$<0.2~\mathrm{km}$), while WPA and DPA capture near-surface wind speed more accurately and agree with scanning lidar and virtual tower data. Above the surface layer, IPA and WPA perform similarly and deviate marginally from the observations. During the transition period (0000 UTC), there are no significant differences among the methods in terms of wind speed, and all closely follow the observational data.

For the stable regime (0600~UTC) (Fig.~\ref{fig:verticalFullDomain2}), the wavelet-based assimilation method captures the low-level jet and near-surface structure well, matching the observations. The indirect assimilation method shows deviations near the surface and slightly underestimates the low-level jet magnitude. Comparison with the virtual tower data is particularly useful below 100 m, where scanning lidar measurements at site A2 are unavailable. The LES wind speed profiles agree well with the virtual tower data during the unstable (Fig.~\ref{fig:verticalFullDomain1}b) and stable (Fig.~\ref{fig:verticalFullDomain2}) regimes, when the virtual tower measurements align with the scanning lidar data. However, during the transition period (Fig.~\ref{fig:verticalFullDomain1}b), a significant discrepancy appears. This mismatch is attributed to the spatial offset between the measurement sites, as the virtual tower data are collected at site A1, whereas the scanning lidar data are obtained from site A2.
 
In terms of wind direction, all methods (direct, indirect, and wavelet) slightly overestimate values near the surface at 1800~UTC (Fig.~\ref{fig:verticalFullDomain1}a). During the transition phase (0000~UTC, Fig.~\ref{fig:verticalFullDomain1}b), the profiles are similar across all methods and match the observations well, with only minor differences above the boundary layer. At 0600~UTC (Fig.~\ref{fig:verticalFullDomain2}), the indirect assimilation method shows deviations throughout the boundary layer, with the largest differences in the surface layer, while the direct and wavelet methods follow the observed wind direction closely.


\begin{figure}[htb!]
    \centering
    \large
    \begin{tabular}{c}
       {\large (a) 2023-09-01 18:00 UTC} (convective)\\
        \includegraphics[width=0.75\linewidth]{Plots/fullDomain/verticalProfiles_1_DPA.png} \\ 
        {\large (b) 2023-09-02 00:00 UTC} (transitional)\\
        \includegraphics[width=0.75\linewidth]{Plots/fullDomain/verticalProfiles_2_DPA.png} \\  
    \end{tabular}
\caption{
Vertical profiles of horizontal wind speed, wind direction, virtual potential temperature, and turbulent kinetic energy from the microscale solver. The microscale solver data are planar-averaged and 1-hour temporally averaged. Panel~(a) corresponds to the convective period (1800~UTC), and panel~(b) corresponds to the transition period (0000~UTC). These vertical profiles are obtained using classical DPA, IPA, and the new WPA profile assimilation method.
}

    \label{fig:verticalFullDomain1}
\end{figure}

\begin{figure}[htb!]
    \centering
    \large
    \begin{tabular}{c}
         {\large 2023-09-02 06:00 UTC} (stable)\\
        \includegraphics[width=0.75\linewidth]{Plots/fullDomain/verticalProfiles_3_DPA.png} \\ 
    \end{tabular}
\caption{
Vertical profiles of horizontal wind speed, wind direction, virtual potential temperature, and turbulent kinetic energy from the microscale solver. The microscale solver data are planar-averaged and 1-hour temporally averaged, and correspond to the stable period (0600~UTC). These vertical profiles are obtained using classical DPA, IPA, and the new WPA profile assimilation method.
}

    \label{fig:verticalFullDomain2}
\end{figure}

For virtual potential temperature during the convective period at 1800~UTC (Fig.~\ref{fig:verticalFullDomain1}a), both indirect and direct approaches reproduce the observed vertical structure well, while WPA shows a deviation in the region above the boundary layer. During the transition period (0000~UTC, Fig.~\ref{fig:verticalFullDomain1}b), the indirect method deviates within and above the boundary layer, whereas the wavelet method remains closer to the observational data. In the stable regime (0600~UTC, Fig.~\ref{fig:verticalFullDomain2}), all methods capture the vertical profiles of virtual potential temperature within the boundary layer, though IPA shows a slight deviation above the ABL.

Finally, for turbulent kinetic energy, the vertical profiles at 1800~UTC, 0000~UTC, and 0600~UTC (Figs.~\ref{fig:verticalFullDomain1} and \ref{fig:verticalFullDomain2}) show distinct patterns. As expected, DPA significantly overestimates TKE throughout the boundary layer, though it converges toward the observations above it. IPA represents the average TKE well within the boundary layer, while WPA lies between DPA and IPA and appears to slightly overpredict the TKE values in the convective and transition regions. 

\subsection{\textbf{Hybrid profile assimilation}}
\label{sec:hybridPA}
\subsubsection{Diurnal evolution of data driven LES}
This section evaluates HGWPA techniques over the full diurnal cycle of the data-driven LES. 
The analysis investigates how the hybrid geostrophic-wavelet approaches (HGWPA1, HGWPA2, and HGWPA3) reproduce the temporal evolution of mean and turbulent flow quantities at multiple heights.
Figures~\ref{fig:timehistoryMeanHGWPA100m}a and~\ref{fig:timehistoryMeanHGWPA100m}b present the time history of mean-flow quantities at $100~\mathrm{m}$ obtained from the hybrid geostrophic--wavelet profile assimilation schemes. For comparison, results from the wavelet-based profile assimilation are shown together with observational and virtual tower data. Figure~\ref{fig:timehistoryMeanHGWPA100m}a compares the horizontal wind speed among the hybrid schemes. HGWPA1 exhibits a noticeable drift relative to both the mesoscale and virtual tower data, whereas HGWPA2 and HGWPA3 do not show such drift but tend to overpredict wind speed during the stable regime, as highlighted in the zoomed-in interval between 0900 and 1500~UTC.
The drift observed in HGWPA1 arises from the forcing implementation. Above the boundary layer, the forcing comprises the inertial and Coriolis terms applied to the mesoscale geostrophic wind. Within the boundary layer, however, the applied forcing remains constant with height, and the vertical gradient of turbulent momentum flux is neglected. The omission of this term leads to a phase-lag in the simulated wind speed relative to the observational data as the flow changes slowly to the applied pressure gradient term. 


Figure~\ref{fig:timehistoryMeanHGWPA100m}b compares wind direction and virtual potential temperature. In HGWPA1, the absence of the turbulent drag term also impacts the estimation of wind direction, which diverges relative to the mesoscale results. HGWPA2 and HGWPA3 reduce this deviation as both schemes implement a turbulent drag force in their algorithm, with HGWPA2 providing the closest agreement. All schemes reproduce the virtual potential temperature accurately, since it is assimilated directly through the wavelet-based profile method.

\begin{figure}[htb!]
    \centering
    \begin{tabular}{c}
        {\large (a)} \\ [2pt]
         \includegraphics[width=0.6\linewidth]{Plots/fullDomain/windSpeed_HGWPA_100m.png} \\
        {\large (b)} \\
        \includegraphics[width=0.6\linewidth]{Plots/fullDomain/windDirection_HGWPA_100m.png} \\
    \end{tabular}
\caption{
Temporal variation of mean flow quantities at $100~\mathrm{m}$ height for the hybrid geostrophic-wavelet and WPA schemes. Mean profiles in panel~(a) include wind speed and a magnified wind speed segment, and panel~(b) includes wind direction and virtual potential temperature.
}
\label{fig:timehistoryMeanHGWPA100m}
\end{figure}

In terms of turbulent flow quantities, all methods capture realistic temporal variations of turbulent kinetic energy (TKE) and shear stress at $100~\mathrm{m}$ (Fig.~\ref{fig:timehistoryTurbHGWPA100m}), reproducing the overall trends observed in the measurements. During the stable regime (0600--1200~UTC on the first day and 0000--1600~UTC on the second day), HGWPA2 exhibits an overestimation of TKE, whereas HGWPA3 performs better owing to its smooth tapering of the frictional drag above the ABL.
The significant overprediction in HGWPA2 on the second day likely stems from the residual layer---a remnant of the previous day’s mixed layer that persists aloft during stable conditions. This layer often retains weak yet sustained turbulence, which can artificially enhance the drag term if not adequately damped. HGWPA3 mitigates this issue by explicitly identifying the ABL height and reducing the drag to zero above it, while HGWPA2 applies drag uniformly without defining the ABL top.
For shear stress (Fig.~\ref{fig:timehistoryTurbHGWPA100m}b), HGWPA1 underperforms in the convective boundary layer because its forcing remains constant with height and fails to respond to evolving boundary-layer structures. HGWPA2 and HGWPA3 capture some degree of coherence in the shear stress field, though less effectively than WPA.

\begin{figure}[htb!]
    \centering
    \begin{tabular}{c}
        {\large (a) } \\ [2pt]
         \includegraphics[width=0.6\linewidth]{Plots/fullDomain/kineticEnergy_HGWPA_100m.png} \\
       {\large (b)}\\
        \includegraphics[width=0.6\linewidth]{Plots/fullDomain/shearStress_HGWPA_100m.png} \\
    \end{tabular}
\caption{
Temporal evolution of turbulent flow quantities at $100~\mathrm{m}$ height using hybrid geostrophic-wavelet schemes and WPA. Panel~(a) includes turbulent kinetic energy, a close-up of turbulent kinetic energy ($1500$--$2100~\mathrm{UTC}$), and panel~(b) shows shear stress, and its detailed section ($1500$--$2100~\mathrm{UTC}$).
}

\label{fig:timehistoryTurbHGWPA100m}
\end{figure}

Figures~\ref{fig:timehistoryMeanHGWPA200m}a and \ref{fig:timehistoryMeanHGWPA200m}b compare the HGWPA methods at $200~\mathrm{m}$. HGWPA2 and HGWPA3 predict mean flow quantities more accurately than at $100~\mathrm{m}$, showing improved agreement with the observational data. However, HGWPA1 continues to exhibit model drift in both wind speed and wind direction. All methods now produce comparable levels of turbulent kinetic energy (TKE) and follow the diurnal trends of the observational results (Figures~\ref{fig:timehistoryTurbHGWPA200m}).

\begin{figure}[htb!]
    \centering
    \begin{tabular}{c}
        {\large (a)}  \\ [2pt]
         \includegraphics[width=0.6\linewidth]{Plots/fullDomain/windSpeed_HGWPA_200m.png} \\
        {\large (b)}  \\
        \includegraphics[width=0.6\linewidth]{Plots/fullDomain/windDirection_HGWPA_200m.png} \\
    \end{tabular}
    \caption{Temporal variation of mean flow quantities at a height of $200~\mathrm{m}$. Mean profiles (a) include wind speed, a magnified wind speed segment, (b) wind direction, and virtual potential temperature.}
    \label{fig:timehistoryMeanHGWPA200m}
\end{figure}

\begin{figure}[htb!]
    \centering
    \begin{tabular}{c}
        {\large (a) }\\ [2pt]
         \includegraphics[width=0.6\linewidth]{Plots/fullDomain/kineticEnergy_HGWPA_200m.png} \\
       {\large (b)}\\
        \includegraphics[width=0.6\linewidth]{Plots/fullDomain/shearStress_HGWPA_200m.png} \\
    \end{tabular}
\caption{Temporal evolution of turbulent flow quantities at a height of $200~\mathrm{m}$. Panel (a) includes turbulent kinetic energy, a close-up of turbulent kinetic energy ($1500$--$2100~\mathrm{UTC}$), and panel~(b) shows shear stress, and its detailed section ($1500$--$2100~\mathrm{UTC}$).}
    \label{fig:timehistoryTurbHGWPA200m}
\end{figure}

\subsubsection{Vertical structure of the atmospheric boundary layer}

Figures~\ref{fig:verticalHGWPAA}a and \ref{fig:verticalHGWPAA}b present vertical profiles of wind speed, wind direction, virtual potential temperature, and turbulent kinetic energy (TKE) for the convective (1800~UTC) and transitional (0000~UTC) periods. Within the convective boundary layer, all schemes reproduce the wind speed profiles reasonably well. For wind direction, HGWPA1 deviates from observations within the ABL but aligns with the other schemes above it, a discrepancy attributable to the absence of the turbulence drag term in HGWPA1. Inclusion of this drag in HGWPA2 and HGWPA3 improves wind direction predictions, yielding results consistent with WPA and mesoscale data.
For TKE, both HGWPA2 and HGWPA3 predict higher values near the surface compared to HGWPA1, WPA, and the observational data. 

\begin{figure}[htb!]
    \centering
    \large
    \begin{tabular}{c}
       {\large (a) 2023-09-01 18:00 UTC} (convective)\\
        \includegraphics[width=0.75\linewidth]{Plots/fullDomain/verticalProfiles_1_HGWPA.png} \\ 
        {\large (b) 2023-09-02 00:00 UTC} (transitional)\\
        \includegraphics[width=0.75\linewidth]{Plots/fullDomain/verticalProfiles_2_HGWPA.png} \\ 
    \end{tabular}
\caption{
Vertical profiles of horizontal wind speed, wind direction, virtual potential temperature, and turbulent kinetic energy from the microscale solver. The microscale solver data are planar-averaged and 1-hour temporally averaged. Panel~(a) corresponds to the convective period (1800~UTC), and panel~(b) corresponds to the transition period (0000~UTC).
}
\label{fig:verticalHGWPAA}
\end{figure}
In the transitional regime, HGWPA2 and HGWPA3 predict higher near-surface wind speeds. However, HGWPA2 gradually converges to the mesoscale and WPA profiles, whereas HGWPA3 maintains higher wind speeds throughout and above the ABL. The discrepancy in HGWPA3 can be attributed to its inaccurate prediction of the ABL height as the flow transitions from an unstable to a stable regime. HGWPA1 underpredicts wind speed in the lower half of the ABL and overpredicts it in the upper half.
For wind direction, WPA and HGWPA2 closely match observations, while HGWPA3 slightly overpredicts it, and HGWPA1 exhibits minimal variation with height within the ABL but aligns near the ABL top. All schemes produce virtual potential temperature profiles that slightly underestimate values in the lower ABL. The TKE profiles display similar trends to the convective case: HGWPA1 aligns more closely with observations, whereas the other schemes predict higher values near the surface, converging with observations further above the ground.

In the stable atmospheric regime at 0600 UTC (Fig.~\ref{fig:verticalHGWPAB}a), the HGWPA2 configuration accurately resolves the low-level jet and aligns with measurements from scanning lidar and WPA data, while, HGWPA1 underestimates the jet's peak speed. Due to the shallow boundary layer height, HGWPA1 yields wind direction predictions that agree with observations across most of the vertical domain extent, except within the ABL closer to the surface. Other methods are able to match the wind direction correctly, exhibiting only minor discrepancies near the surface. Regarding turbulent kinetic energy (TKE), HGWPA2 overestimates values throughout the ABL, attributable to the influence of the residual layer, which generates pronounced gradients in turbulent shear stress. HGWPA1 and HGWPA3, however, produce TKE profiles that more closely match the observed data.

An overall comparison of the three hybrid strategies reveals distinct performance characteristics across stability regimes. HGWPA1, which omits the frictional drag term in large-scale pressure gradient computations, underperforms in unstable and neutral conditions due to the larger boundary layer height and the assumption of a constant pressure gradient within the ABL. This assumption, common in canonical studies, does not hold in realistic MMC simulations. HGWPA2, which includes the frictional drag term, performs well in unstable and stable regimes. However, it overestimates TKE in the stable regime due to the residual layer aloft, where turbulence intensity is weak but higher than in the underlying SBL. This intensity varies significantly with height, leading to large vertical gradients in turbulent shear stress. HGWPA3, which tapers the frictional drag term to zero above the ABL, performs better in the stable regime but worse than HGWPA2 in unstable and transition conditions.

\begin{figure}[htb!]
    \centering
    \large
    \begin{tabular}{c}
        {\large 2023-09-02 06:00 UTC} (stable)\\
        \includegraphics[width=0.75\linewidth]{Plots/fullDomain/verticalProfiles_3_HGWPA.png} \\ 
    \end{tabular}
\caption{
Vertical profiles of horizontal wind speed, wind direction, virtual potential temperature, and turbulent kinetic energy from the microscale solver. The microscale solver data are planar-averaged and 1-hour temporally averaged, and correspond to the stable period (0600~UTC).
}
\label{fig:verticalHGWPAB}
\end{figure}

\section{Discussion}
\label{sec:discussion}
\subsection{Performance of assimilation methods}
While one-hour average vertical plots or 10-minute average time series data at specific heights in the Sect. \ref{sec:results} clearly indicate superior performance for certain models, it remains challenging to definitively determine which model performs best over the full diurnal cycle or across different stability regimes solely from visual inspection. Therefore, here we quantitatively assess the performance of the six MMC models using the root mean square error (RMSE) and Willmott’s index of agreement ($d_{\mathcal{V}}$) \citep{willmott1981validation} to rank them based on prediction accuracy and variability. We use time-height data, comparing microscale outputs from each model with observational data over a 24-hour period from 0600 UTC on September 2, 2023, to 0600 UTC on September 3, 2023. All datasets are averaged over ten-minute intervals, considering only data from 100~m to 1100~m. For each time and height, we form pair-wise values of model predictions and observations for all variables, including primitive variables (wind speed, direction, virtual potential temperature) and derived variables (turbulent kinetic energy, shear stress). All variables are normalized to the range [0, 1] using min-max normalization to ensure consistent scaling across metrics.

For a given variable $\mathcal{V}$, let $x_{i,j}^{(\mathcal{V})}$ represent the observational value at time index $i$ and vertical level $j$, and $y_{i,j}^{(\mathcal{V})}$ represent the corresponding model prediction, both after min-max normalization. The root mean square error (RMSE) for variable $\mathcal{V}$ is computed across all pair-wise times and heights as:
\begin{equation}
    \text{RMSE}_\mathcal{V} = \sqrt{\frac{1}{N_\mathcal{V}} \sum_{i,j} \big(y_{i,j}^{(\mathcal{V})} - x_{i,j}^{(\mathcal{V})}\big)^2},
    \label{eq:rmse}
\end{equation}
where $N_\mathcal{V}$ is the total number of valid paired points for variable $\mathcal{V}$.

In addition to the root mean square error (RMSE), which quantifies the average error between model predictions and observations, we compute Willmott’s index of agreement ($d_\mathcal{V}$) \citep{willmott1981validation} for each variable to measure how well predictions capture the variability of observed data. This index assesses the relative covariability of predictions and observations, evaluating how closely their deviations correspond around the mean of the observations, which acts as a reference for the expected true value. By combining RMSE’s measure of prediction error with $d_\mathcal{V}$’s evaluation of variability patterns, we rank models based on both their accuracy and their ability to reflect observed data trends. $d_\mathcal{V}$ is computed as
\begin{equation}
    d_\mathcal{V} = 1 - \frac{\sum_{i,j} (y_{i,j}^{(\mathcal{V})} - x_{i,j}^{(\mathcal{V})})^2}{\sum_{i,j} \big(|y_{i,j}^{(\mathcal{V})} - \bar{x}^{(\mathcal{V})}| + |x_{i,j}^{(\mathcal{V})} - \bar{x}^{(\mathcal{V})}|\big)^2},
    \label{eq:willmott}
\end{equation}
where $\bar{x}^{(\mathcal{V})}$ is the mean of the observations for variable $\mathcal{V}$. The index $d_\mathcal{V}$ ranges from 0 to 1, with values closer to 1 indicating better agreement.

Finally, we calculate a composite score for each model, denoted $\mathcal{S}_c$, by combining RMSE and the Willmott’s index ($d_\mathcal{V}$) across all variables with assigned weights:
\begin{equation}
\mathcal{S}_\mathrm{c} = \sum_\mathcal{V} \mathcal{W}_\mathcal{V} \Big[ \alpha \big(1 - \text{RMSE}_\mathcal{V}\big) + (1 - \alpha) \, d_\mathcal{V} \Big],
\label{eq:compositescore}
\end{equation}
where $\mathcal{W}_\mathcal{V}$ is the weight assigned to variable $\mathcal{V}$, and $\alpha = 0.5$. The weighting scheme reflects the relative importance of each variable in evaluating model performance. Lower weights are assigned to primitive variables--wind speed ($\mathcal{W}_{\mathrm{mag}} = 0.15$), wind direction ($\mathcal{W}_{\mathrm{dir}} = 0.15$), and virtual potential temperature ($\mathcal{W}_{\theta} = 0.15$)---while higher weights are assigned to derived quantities such as turbulent kinetic energy ($\mathcal{W}_{\mathrm{tke}} = 0.30$) and shear stress ($\mathcal{W}_{\tau} = 0.35$). This approach emphasizes the sensitivity of turbulence-related quantities while maintaining balance across all evaluated parameters. The resulting score enables ranking of model performance over the full diurnal cycle and within distinct stability regimes of the 24-hour observational period. In this analysis, we consider analyzing models performance for full diurnal cycle (24 hour period), convective period (5 hour period from 1700 to 2200 UTC), and stable period (5 hour period from 0700 to 1200 UTC).

\begin{figure}[htp!]
    \centering
    \includegraphics[width=1.0\linewidth]{Plots/stat_diurnal.png}
    \caption{Scatter plots comparing model data against observed data for the full diurnal cycle over a 24-hour period across six MMC models (columns from left to right: WPA, HGWPA2, HGWPA3, HGWPA1, IPA, DPA). Rows represent different primitive and derived variables, including wind speed ($S$) (top row), direction (Dir) (second row), virtual potential temperature ($\theta$) (third row), turbulent kinetic energy ($k$) (fourth row), and shear-stress ($\tau_z$) (bottom row). Each plot includes a y = x reference line (black), a 95\% confidence ellipse (shaded region), and the Willmott's index of agreement ($d_\mathcal{V}$) value indicating model performance.}
    \label{fig:scatterplot}
\end{figure}

Table~\ref{tab:fulldiurnal} presents the normalized RMSE and Willmott index $d_\mathcal{V}$ for primitive and derived variables across six MCMC models over the full diurnal cycle (a 24-hour period capturing day-night variations in atmospheric conditions). These metrics evaluate how well the models (DPA, IPA, WPA, HGWPA1, HGWPA2, and HGWPA3) predict primitive and derived variables. Across all models except HGWPA1, mean flow variables ($S$, Dir, $\theta_v$) are predicted with high accuracy over the full cycle, showing low RMSE (mostly $<$0.05) and $d_\mathcal{V}$ values near 1 (indicating strong agreement with observations). For instance, wind speed $S$ has RMSE below $0.04$ for DPA, IPA, WPA, and HGWPA2, with $d > 0.99$, reflecting robust capture of diurnal wind patterns. Virtual potential temperature $\theta_v$ shows similar performance, with DPA achieving the lowest RMSE (0.0359), followed by WPA and HGWPA3 ($0.0241$ and $0.0253$). Wind direction is generally well-reproduced, though HGWPA1 shows higher RMSE ($0.2555$) and lower $d_\mathcal{V}$ (0.9116). The higher RMSE and lower $d_{\mathcal{V}}$ for HGWPA1 in both wind speed and direction are consistent with the model drift discussed earlier (Figs.~\ref{fig:timehistoryMeanHGWPA100m} and~\ref{fig:timehistoryMeanHGWPA200m}) and are further evident in the scatter plots (Fig.~\ref{fig:scatterplot}, column~4, rows~1–2), where HGWPA1 shows larger deviations from the observational data.

\begin{table}[htb!]
\centering
\caption{Normalized root mean square error (RMSE, top panel) computed using Eq.~\ref{eq:rmse}, and Willmott index ($d$, bottom panel) computed using Eq.~\ref{eq:willmott}, for the full diurnal cycle. Results are shown for six different MMC models (DPA, IPA, WPA, HGWPA1, HGWPA2, and HGWPA3), and for variables wind speed ($S$), wind direction (Dir), virtual potential temperature ($\theta_v$), turbulent kinetic energy ($k$), and vertical shear stress ($\tau_z$).}
\label{tab:fulldiurnal}
\begin{tabular}{lcccccc}
\toprule
Variable & DPA & IPA & WPA & HGWPA1 & HGWPA2 & HGWPA3 \\
\midrule
\multicolumn{7}{c}{\textbf{RMSE}} \\
\midrule
$S$        & 0.0167 & 0.0399 & 0.0363 & 0.1223 & 0.0309 & 0.0704 \\
Dir        & 0.0357 & 0.0548 & 0.0497 & 0.2555 & 0.0618 & 0.1254 \\
$\theta_v$ & 0.0079 & 0.0227 & 0.0241 & 0.0174 & 0.0378 & 0.0253 \\
$k$        & 0.5270 & 0.1879 & 0.1802 & 0.1536 & 0.1809 & 0.1511 \\
$\tau_z$   & 0.1807 & 0.1275 & 0.1154 & 0.1219 & 0.1186 & 0.1192 \\
\midrule
\multicolumn{7}{c}{\textbf{Willmott index ($d_\mathcal{V}$)}} \\
\midrule
$S$        & 0.9987 & 0.9925 & 0.9939 & 0.9108 & 0.9955 & 0.9719 \\
Dir        & 0.9983 & 0.9960 & 0.9967 & 0.9116 & 0.9948 & 0.9768 \\
$\theta_v$ & 0.9996 & 0.9964 & 0.9959 & 0.9979 & 0.9901 & 0.9955 \\
$k$        & 0.5695 & 0.6932 & 0.8378 & 0.8645 & 0.8232 & 0.8652 \\
$\tau_z$   & 0.5587 & 0.4849 & 0.6194 & 0.4873 & 0.5817 & 0.6015 \\
\bottomrule
\end{tabular}
\end{table}

Turbulence-related variables ($k$ and $\tau_z$) exhibit greater variability among the models, highlighting the inherent difficulty of reproducing turbulent processes across the diurnal cycle. For $k$, DPA performs the worst (RMSE = 0.5270, $d_{\mathcal{V}} = 0.5695$), while the hybrid models HGWPA1 and HGWPA3 achieve substantially better accuracy (RMSE $\approx$ 0.15, $d_{\mathcal{V}} \approx$ 0.86), indicating that the wavelet-based blending improves turbulence prediction. Similarly, for $\tau_z$, DPA again underperforms (RMSE = 0.1807, $d_{\mathcal{V}} = 0.5587$), whereas WPA attains the lowest RMSE (0.1154) and highest $d_{\mathcal{V}}$ (0.6194), outperforming the hybrid models for this metric.

Table~\ref{tab:convective} summarizes the normalized RMSE and Willmott index $d_\mathcal{V}$ for the convective period, showing model performance patterns that generally align with the full diurnal cycle results in Table~\ref{tab:fulldiurnal}. Mean flow variables ($S$, Dir, $\theta_v$) are well captured across most models, with low RMSE ($<0.06$ for $S$ and $\theta_v$, $<0.08$ for Dir) and $d_\mathcal{V} > 0.98$ for DPA, IPA, WPA, and HGWPA2. For example, DPA achieves the lowest RMSE for $S$ (0.0353) and $\theta_v$ (0.0230) with $d_\mathcal{V} > 0.99$, while WPA leads in Dir ($d_\mathcal{V}=0.9905$). HGWPA1 remains a clear outlier, with much higher RMSE for $S$ (0.2767) and Dir (0.4091) and lower $d_\mathcal{V}$ values (0.5364 and 0.7095), reflecting the same drift seen in the full cycle. HGWPA3 also shows elevated RMSE for $S$ (0.2046, $d_\mathcal{V}=0.7658$) and Dir (0.1484, $d_\mathcal{V}=0.9477$) compared to the non-hybrid models and HGWPA2, but still performs better than HGWPA1.

Turbulence variables ($k$ and $\tau_z$) vary more under convective conditions. DPA performs poorly in $k$ (RMSE=0.8192, $d_\mathcal{V}=0.2801$) but reasonably in $\tau_z$ (RMSE=0.2312, $d_\mathcal{V}=0.3640$). WPA again excels in $\tau_z$ (RMSE=0.2311, $d_\mathcal{V}=0.4086$), consistent with its full-cycle performance, while the hybrid models (HGWPA1, HGWPA2, HGWPA3) retain an advantage in $k$ prediction (RMSE $\approx$0.24–0.27, $d_\mathcal{V} \approx 0.51$), outperforming DPA and IPA (RMSE=0.3720, $d_\mathcal{V}=0.4182$). WPA achieves the lowest $k$ RMSE overall (0.2558, $d_\mathcal{V}=0.5374$). Similar performance trends are also observed during the stable period, as shown in Table~\ref{tab:stable}.


% Convective Category Tables
\begin{table}[htb!]
\centering
\caption{Normalized root mean square error (RMSE, top panel) computed using Eq.~\ref{eq:rmse}, and Willmott index ($d$, bottom panel) computed using Eq.~\ref{eq:willmott}, for the convective period. Results are shown for six different MMC models (DPA, IPA, WPA, HGWPA1, HGWPA2, and HGWPA3), and for wind speed ($S$), wind direction (Dir), virtual potential temperature ($\theta_v$), turbulent kinetic energy ($k$), and vertical shear stress ($\tau_z$).}
\label{tab:convective}
\begin{tabular}{lcccccc}
\toprule
Variable & DPA & IPA & WPA & HGWPA1 & HGWPA2 & HGWPA3 \\
\midrule
\multicolumn{7}{c}{\textbf{RMSE}} \\
\midrule
$S$        & 0.0353 & 0.0496 & 0.0466 & 0.2767 & 0.0598 & 0.2046 \\
Dir        & 0.0639 & 0.0736 & 0.0718 & 0.4091 & 0.1041 & 0.1484 \\
$\theta_v$ & 0.0230 & 0.0535 & 0.0470 & 0.0494 & 0.0474 & 0.0466 \\
$k$        & 0.8192 & 0.3720 & 0.2558 & 0.2444 & 0.2682 & 0.2713 \\
$\tau_z$   & 0.2312 & 0.2631 & 0.2311 & 0.2510 & 0.2353 & 0.2412 \\
\midrule
\multicolumn{7}{c}{\textbf{Willmott index ($d_\mathcal{V}$)}} \\
\midrule
$S$        & 0.9920 & 0.9842 & 0.9861 & 0.5364 & 0.9761 & 0.7658 \\
Dir        & 0.9925 & 0.9899 & 0.9905 & 0.7095 & 0.9794 & 0.9477 \\
$\theta_v$ & 0.9976 & 0.9863 & 0.9894 & 0.9885 & 0.9894 & 0.9898 \\
$k$        & 0.2801 & 0.4182 & 0.5374 & 0.5082 & 0.5397 & 0.5283 \\
$\tau_z$   & 0.3640 & 0.4282 & 0.4086 & 0.4185 & 0.3826 & 0.3629 \\
\bottomrule
\end{tabular}
\end{table}



% Stable Category Tables
\begin{table}[htb!]
\centering
\caption{Normalized root mean 
square error (RMSE, top panel) computed using Eq.~\ref{eq:rmse}, and Willmott index ($d$, bottom panel) computed using Eq.~\ref{eq:willmott}, for the stable period. Results are shown for six different MMC models (DPA, IPA, WPA, HGWPA1, HGWPA2, and HGWPA3), and for wind speed ($S$), wind direction (Dir), virtual potential temperature ($\theta_v$), turbulent kinetic energy ($k$), and vertical shear stress ($\tau_z$).}
\label{tab:stable}
\begin{tabular}{lcccccc}
\toprule
Variable & DPA & IPA & WPA & HGWPA1 & HGWPA2 & HGWPA3 \\
\midrule
\multicolumn{7}{c}{\textbf{RMSE}} \\
\midrule
$S$        & 0.0152 & 0.0321 & 0.0304 & 0.2275 & 0.0348 & 0.1033 \\
Dir        & 0.0388 & 0.0872 & 0.0675 & 0.4089 & 0.0799 & 0.2588 \\
$\theta_v$ & 0.0100 & 0.0414 & 0.0222 & 0.0143 & 0.0517 & 0.0331 \\
$k$        & 0.3982 & 0.1051 & 0.1045 & 0.0981 & 0.3319 & 0.1036 \\
$\tau_z$   & 0.1279 & 0.0411 & 0.0411 & 0.0409 & 0.1134 & 0.0410 \\
\midrule
\multicolumn{7}{c}{\textbf{Willmott index ($d_\mathcal{V}$)}} \\
\midrule
$S$        & 0.9989 & 0.9952 & 0.9957 & 0.6436 & 0.9942 & 0.9304 \\
Dir        & 0.9991 & 0.9952 & 0.9971 & 0.9054 & 0.9959 & 0.9593 \\
$\theta_v$ & 0.9992 & 0.9868 & 0.9962 & 0.9984 & 0.9801 & 0.9916 \\
$k$        & 0.0290 & 0.3807 & 0.3592 & 0.3713 & 0.0647 & 0.3400 \\
$\tau_z$   & 0.0518 & 0.2023 & 0.1933 & 0.1985 & 0.0383 & 0.1770 \\
\bottomrule
\end{tabular}
\end{table}

Finally, Table~\ref{tab:finalscore} presents the final composite score computed using Eq.~\ref{eq:compositescore} for the full diurnal cycle, convective period, and stable period. The composite score shows that wavelet-profile assimilation performs best in all three cases, with IPA and WPA nearly equal in the stable period. HGWPA2 ranks second overall, performing very well during the full diurnal cycle and convective period but showing weak performance in the stable period. This limitation is consistent with the vertical profile of turbulent kinetic energy, where HGWPA2 produces elevated TKE (Fig.~\ref{fig:verticalHGWPAB}). This problem of elevated TKE is alleviated when ABL height information is included, as seen in HGWPA3, which performs better than HGWPA2.

\begin{table}[htbp]
\centering
\caption{Composite scores $\mathcal{S}_c$ for six MMC models. Maximum values in each column are highlighted.}
\label{tab:finalscore}
\begin{tabular}{lccc}
\toprule
Model & Full diurnal cycle & Convective period & Stable period\\
\midrule
WPA & \cellcolor{gray!30}0.9526 & \cellcolor{gray!30}0.8333 & 0.8300\\
HGWPA2 & 0.9411 & 0.8215 & 0.7071\\
HGWPA3 & 0.9458 & 0.7826 & 0.7957\\
HGWPA1 & 0.9025 & 0.7313 & 0.7608\\
IPA & 0.9034 & 0.7946 & \cellcolor{gray!30}0.8308\\
DPA & 0.8427 & 0.7069 & 0.7013\\
\bottomrule
\end{tabular}
\end{table}

\subsection{\textbf{Turbulence characterization}}
To assess the fidelity of meso-to-microscale coupling, the resolved velocity signal, $\overline{\bm{u}}(t)$, is sampled from a probe located on the vertical midplane of the computational domain at a height of $z = 110~\mathrm{m}$. The sampling frequency of the resolved velocity $\overline{\bm{u}}(\bm x, t)$ is $f_s = 1~\mathrm{Hz}$. We have also collected the high-frequency mesoscale data, which was sampled at a frequency of $f_s = 0.25~\mathrm{Hz}$. These velocity signals are analyzed to examine the spectral energy distribution of eddying motions and to evaluate how effectively different profile assimilation techniques reproduce the inertial sub-range and dissipation characteristics within the MMC framework. 

The spectral analysis is performed using the horizontal wind speed
\begin{equation}
S(t) = \sqrt{U^2(t) + V^2(t)},
\end{equation}
where $U(t)$ and $V(t)$ are the streamwise and spanwise components, respectively. The fluctuating signal is obtained as
\begin{equation}
S'(t) = S(t) - \overline{S},
\end{equation}
where $\overline{S}$ is the temporal mean over the analysis period.

The one-dimensional energy spectra are computed from $1~\mathrm{hr}$ of data centered around 1800~UTC, representing the convective boundary layer, and 0600~UTC, representing the stable boundary layer, using the Welch’s method \citep{welch2003use}, with $10~\mathrm{min}$ segments and a Hanning window function. This method was also employed in the previous MMC study by \citet{allaerts2020development}. Let $S'(t)$ denote the fluctuating component of the microscale or mesoscale velocity signal. In the classical periodogram approach, the power spectral density (PSD) estimate $\hat{\mathcal P}_{uu}(f)$ is given by

\begin{equation}
\hat{\mathcal P}_{uu}(f) = \frac{1}{\mathcal{T}_s} \left| \frac{1}{\mathcal N} \sum_{n=0}^{\mathcal N-1} S'(n \, \Delta t) \, e^{-i 2 \pi f n \Delta t} \right|^2,
\end{equation}
where $\mathcal N$ is the total number of samples, $\Delta t$ is the sampling interval, and $\mathcal T_s = \mathcal N \Delta t$ is the signal length which is $1~\mathrm{hr}$ in the present case. Welch’s method improves the variance properties of this estimate by dividing the signal into $\mathcal K$ overlapping segments of length $\mathcal {L}$ samples, applying a window function $\mathcal{W}(n)$ to each segment to reduce spectral leakage, and computing the discrete Fourier transform of each windowed segment to obtain the modified periodogram
\begin{equation}
\hat{\mathcal P}_{uu}^{(m)}(f) = \frac{1}{\mathcal U \, \mathcal L f_s} 
\left| \sum_{n=0}^{\mathcal L - 1} S'(n + m \mathcal D) \, \mathcal{W}(n) \, 
e^{-i 2 \pi f n / f_s} \right|^2,
\end{equation}
where $\mathcal U = \tfrac{1}{\mathcal L} \sum_{n=0}^{\mathcal L-1} \mathcal{W}^2(n)$ 
is the window normalization factor, $\mathcal D$ is the segment shift in samples, 
$n$ is the sample index within each segment, and $m$ is the segment index. 
Averaging the $\mathcal K$ modified periodograms yields
\begin{equation}
\hat{\mathcal P}_{uu}^{\mathrm{Welch}}(f) = \frac{1}{\mathcal K} \sum_{m=0}^{\mathcal K -1} 
\hat{\mathcal P}_{uu}^{(m)}(f).
\end{equation}

Finally, the one-dimensional energy spectrum $E_{uu}(f)$ is obtained from the PSD as
\begin{equation}
E_{uu}(f) = 2 \, \hat{\mathcal P}_{uu}^{\mathrm{Welch}}(f),
\end{equation}
for $f > 0$, where the factor of 2 accounts for conversion from a two-sided to a one-sided spectrum. This procedure yields a statistically converged estimate of the turbulent kinetic energy distribution across temporal frequencies, enabling assessment of the inertial sub-range scaling and dissipation range behavior.

Figures~\ref{fig:euu_dpa}a and \ref{fig:euu_dpa}b present the streamwise velocity energy spectra obtained using DPA, IPA, WPA, and mesoscale data for unstable and stable atmospheric conditions, respectively. In the unstable case (Fig.~\ref{fig:euu_dpa}a), the DPA approach exhibits higher spectral energy across most frequencies compared to the mesoscale spectrum, indicating an overestimation of turbulent kinetic energy (TKE). Nevertheless, DPA reproduces the characteristic $-5/3$ slope within the inertial subrange. In contrast, IPA underestimates the spectral energy, while WPA closely matches the mesoscale spectrum and accurately represents the inertial subrange. Regarding the high-frequency tail, the IPA-based MMC yields the steepest spectral decay, whereas DPA delays the onset of dissipation, with WPA exhibiting an intermediate behavior.

In the stable case (Fig.~\ref{fig:euu_dpa}b), the DPA approach significantly overestimates the spectral energy across all resolved frequencies compared to the mesoscale spectrum. The mesoscale data, in contrast, exhibit a pronounced dip in spectral energy, indicating a lack of smaller-scale motions that are expected to be present within the stable boundary layer. This limitation highlights the benefit of MMC, which allows the microscale solver to recover part of the missing high-frequency content. Both IPA and WPA capture a sufficient portion of the inertial subrange and provide a more realistic spectral behavior relative to the mesoscale data. Overall, IPA and WPA perform comparably well under stable conditions, with each reproducing a physically plausible distribution of energy across scales.

\begin{figure}[htb!]
    \centering
    \large
    \begin{tabular}{cc}
    {\large (a) Unstable}  & {\large (b) Stable} \\
        \includegraphics[width=0.46\linewidth]{Plots/fullDomain/Euu_unstable.png} &
        \includegraphics[width=0.46\linewidth]{Plots/fullDomain/Euu_stable.png} \\ [2pt]
   \end{tabular}
\caption{
Turbulence energy spectra of horizontal wind speed at $110~\mathrm{m}$, computed from 1-hour averages centered at 1800~UTC (unstable) and 0600~UTC (stable). Results are shown for direct profile assimilation (DPA), indirect profile assimilation (IPA), wavelet-based profile assimilation (WPA), and mesoscale observational data; panels show (a) unstable and (b) stable conditions.
}

\label{fig:euu_dpa}
\end{figure}

Next, we examine the turbulent energy spectra obtained using the HGWPA schemes, with the wavelet-based assimilation shown for reference. Figures~\ref{fig:euu_hgpwa}a and \ref{fig:euu_hgpwa}b present the estimated spectral density for unstable and stable conditions, respectively. During the convective period (Fig.~\ref{fig:euu_hgpwa}a), the HGWPA2 scheme performs well, showing close agreement with both the WPA results and the mesoscale spectrum across the inertial subrange. In contrast, HGWPA1 underestimates the spectral energy and exhibits a steeper high-frequency decay, indicative of excessive dissipation. The HGWPA3 scheme, on the other hand, exhibits an attenuated spectral decay at high frequencies, indicating insufficient representation of small-scale dissipation and leading to elevated energy levels in the dissipation range.

For the stable case (Fig.~\ref{fig:euu_hgpwa}b), HGWPA1 does not reproduce any inertial subrange. This behavior can be attributed to the shallow boundary-layer depth under stable stratification and the presence of a residual layer above it. Since HGWPA1 does not include a turbulent drag term and applies uniform forcing below the ABL height, the scheme cannot generate the scale interactions necessary to sustain an inertial subrange. In contrast, HGWPA3 shows good agreement with the observational data spectrum. However, it appears that the observational data spectrum itself does not capture the turbulence scales characteristic of the stable boundary layer, indicating that even though HGWPA3 follows the mesoscale energy distribution, the representation of small-scale turbulence remains limited. All other schemes are able to reproduce the inertial subrange except HGWPA1. Overall, these results highlight the importance of mesoscale--microscale coupling, which incorporates the large-scale variability from mesoscale forcing while retaining the small-scale turbulent structures that are essential for assessing physical phenomena within the atmospheric boundary layer, such as wind-farm performance and wildfire spread.

\begin{figure}[htb!]
    \centering
    \large
    \begin{tabular}{cc}
    {\large (a) Unstable}  & {\large (b) Stable} \\
        \includegraphics[width=0.48\linewidth]{Plots/fullDomain/Euu_unstable_hgwpa.png} 
    &
        \includegraphics[width=0.48\linewidth]{Plots/fullDomain/Euu_stable_hgwpa.png} \\ [2pt]
   \end{tabular}
\caption{
Turbulence energy spectra of horizontal wind speed at $110~\mathrm{m}$, computed from 1-hour averages centered at 1800~UTC (unstable) and 0600~UTC (stable). Results are shown for HGWPA1, HGWPA2, HGWPA3, WPA, and mesoscale observational data; panels show (a) unstable and (b) stable conditions.
}
\label{fig:euu_hgpwa}
\end{figure}

% \clearpage
\section{Conclusion}
\label{sec:conclusions}

This study developed and evaluated novel mesoscale-to-microscale coupling (MMC) approaches for large-eddy simulations, focusing on improving the representation of atmospheric boundary layer (ABL) processes across stability regimes. Two new schemes were introduced: a wavelet-based profile assimilation (WPA) method, which retains the error-based formulation but incorporates a wavelet scale-aware filter to account for local variations in the vertical structure, and a physics-based hybrid geostrophic–wavelet profile assimilation (HGWPA) scheme, which applies geostrophic balance for momentum source terms and wavelet-based profile assimilation for virtual potential temperature. The wavelet filter adaptively retains scale-dependent components of mesoscale tendencies, reducing spurious small-scale oscillations while preserving physically relevant features. The methods were evaluated over a full diurnal cycle using mesoscale forcing derived from observations gathered during the American WAKE experimeNt (AWAKEN) project, and validated against scanning lidar observations, virtual tower diagnostics, and Weather Research and Forecasting (WRF) results.  

Analysis of diurnal time series, vertical profiles, and turbulence energy spectra showed that WPA captured mean-flow quantities, turbulent kinetic energy (TKE), and inertial subrange scaling across all stability regimes while recovering small-scale motions absent in the mesoscale input. The scale-aware wavelet filtering helped maintain vertical coherence and improved the physical consistency of the imposed tendencies without excessive smoothing. The hybrid schemes improved wind direction representation within the ABL through the inclusion of a turbulence drag term in the geostrophic momentum forcing. Among the hybrid methods, HGWPA2 provided the most consistent results across regimes, although it overpredicted TKE near-surface in the stable boundary layer due to the presence of a residual layer above the jet.  

The developed MMC framework, with its scale-aware filtering capability, can be applied to realistic wind farm simulations and other ABL studies requiring accurate mesoscale-to-microscale coupling, including further investigations within the AWAKEN project. Future work will focus on dynamically blending physics-based and error-based forcings in time according to stability conditions, aiming to reduce TKE overprediction in stable regimes. In addition, a machine-learning-based technique that can adaptively select the optimal wavelet filter and levels of decomposition depending on the stability regime, turbulence intensity, or other ABL parameters will be explored.

%\section{Reference Formatting}
\begin{acknowledgements}
This research has been supported by the Natural Sciences and Engineering Research Council of Canada (grant no. 556326) in partnership with UL Renewables. Computational resources provided by the Digital Research Alliance of Canada are gratefully acknowledged.
 
\end{acknowledgements} 

\section*{Code Availability}
The TOSCA code used in this study is available at: \url{https://github.com/sebastipa/TOSCA}

\clearpage
\appendix
\section*{Appendix}
\addcontentsline{toc}{section}{Appendix}
\setcounter{section}{0} % reset section counter
\renewcommand{\thesection}{A\arabic{section}} % ensure A1 format
\section{Grid Resolution Analysis}
\label{sec:GRAstudy}

A grid resolution analysis for a dynamically evolving ABL, where LES was driven by prescribed measurements over a full diurnal cycle, is performed. Assessing the grid resolution requirements in this context is particularly challenging as the ABL is continuously changing while transitioning between the convective and stable regimes. To address this challenge, we examine how different grid configurations influence key physical quantities under varying stability regimes. The grid study included isotropic grids of streamwise, spanwise and vertical sizes $10~\mathrm{m} \times 10~\mathrm{m} \times 10~\mathrm{m}$, $20~\mathrm{m} \times 20~\mathrm{m} \times 20~\mathrm{m}$, $30~\mathrm{m} \times 30~\mathrm{m} \times 30~\mathrm{m}$ and $40~\mathrm{m} \times 40~\mathrm{m} \times 40~\mathrm{m}$, as well as anisotropic grids of size $20~\mathrm{m} \times 20~\mathrm{m} \times 10~\mathrm{m}$, $30~\mathrm{m} \times 30~\mathrm{m} \times 10~\mathrm{m}$ and $40~\mathrm{m} \times 40~\mathrm{m} \times 10~\mathrm{m}$. The domain size is kept constant for all the cases as $3~\mathrm{km} \times 3~\mathrm{km} \times 2~\mathrm{km}$ and flow statistics are computed over a two-hour averaging window corresponding to the convective and stable phases of the diurnal cycle. The convective regime is taken as the daytime period from $1700$ UTC to $1900$ UTC on 1st September $2023$ with an average Obukhov Length  $L = -24.18~\mathrm{m}$, while the stable regime is taken as night-time period from $0500$ UTC to $0700$ UTC on 2nd September $2023$ with an average Obukhov length $L = 15.23~\mathrm{m}$. All the cases are performed using DPA method to drive the LES. 

Various diagnostic metrics are used to assess the grid resolution requirements: (i) vertical profiles of horizontal wind speed ($S$) and squared friction velocity ($u_{*}^2$), which characterize the resolved mean flow and provide a surface--stress-related measure; (ii) turbulence-resolving diagnostics, including the resolved kinetic energy ($k_{res}$), the LES effective resolution index ($\gamma$), and energy spectra; and (iii) two-point correlation functions to assess the spatial coherence and structure of the flow. All vertical profiles are obtained by first applying horizontal plane averaging and then performing temporal averaging over a two-hour time window. This approach is consistent with resolution assessment frameworks established in previous work by \citet{pope2001turbulent}, \citet{davidson2009large}, \citet{chow2003further}, and \citet{wurps2020grid}.

Figures~\ref{fig:GRAUmag}a and~\ref{fig:GRAUmag}b present the horizontal wind speed $S$ as a function of height ($z$), normalized by the domain height $H$, for the convective and stable regimes, respectively. The results show that all grid resolutions produce similar mean profiles in both regimes. In the stable regime, all grids are able to capture the low-level jet structure, indicating consistency in reproducing mesoscale features of the flow. The convective regime exhibits a nearly uniform profile within the boundary layer due to strong turbulent mixing. However, within the surface layer, finer isotropic and anisotropic resolution cases converge to a single wind speed profile, while coarser isotropic cases ($30\times30\times30$ and $40\times40\times40$) tend to overpredict the wind speed. The coarse anisotropic grid ($40\times40\times10$) underpredicts the wind speed near the surface, likely due to the large aspect ratio of the mesh.

\begin{figure}[htb!]
\centering
\begin{tabular}{cc}
{\large (a)} & {\large (b)} \\
\includegraphics[width=0.35\textwidth]{Plots/gridStudy/Plots/U_mag_Profile_unstable.png} &
\includegraphics[width=0.35\textwidth]{Plots/gridStudy/Plots/U_mag_Profile_stable.png}
\end{tabular}
\caption{Vertical profiles of horizontal wind-speed obtained from dynamically evolving LES of the diurnal cycle. Profiles were obtained by averaging over two-hour windows for (a) convective ($1700-1900$ UTC) and (b) stable ($0500-0700$ UTC) stratification at various resolutions.}
\label{fig:GRAUmag}
\end{figure}

To quantify surface-layer stress and momentum exchange, we compute the squared friction velocity using $u_*^2 = \left(\overline{u'w'}^{\,2} + \overline{v'w'}^{\,2} \right)^{1/2}$. Fig.~\ref{fig:GRAustar2}a shows the vertical profile of $u_*^2$ for the convective boundary layer with significant variations across the grid resolutions within the surface layer of the ABL, where turbulence activity is most pronounced. Within the surface layer, the finest grid (10$\times$10$\times$10) captures the highest shear stress values. Moreover, the next finer grids 20$\times$20$\times$10 and 20$\times$20$\times$20 predict lower values but converge closely with each other. This indicates that the intermediate grids are resolving the surface-layer turbulence reasonably well, although they do not reproduce the peak magnitudes captured by the finest grid. In Fig.~\ref{fig:GRAustar2}b, presence of the residual layer above the stable boundary layer is captured by the finer vertical grid resolution of 10 m, while they are not seen in the coarser vertical grid resolution of 20 m and higher. It is also noted that the coarser horizontal resolution grids predict higher values of $u_*^2$ within the residual layer. This behavior may be linked to the inability of these grids to adequately capture the smaller-scale turbulent motions that govern entrainment across the inversion layer. \citet{Allaerts_Meyers_2017} demonstrated that the characteristic size of these entraining eddies can be represented by the Ellison scale, which defines the vertical extent of temperature fluctuations relative to the mean stratification. Insufficient resolution near the inversion, particularly in the vertical direction, can therefore prevent the model from resolving the Ellison-scale motions responsible for entrainment, leading to over-energetic predictions of shear production and elevated $u_*^2$ in the residual layer.


\begin{figure}[htb!]
\centering
\begin{tabular}{cc}
{\large (a)} & {\large (b)} \\
\includegraphics[width=0.35\textwidth]{Plots/gridStudy/Plots/Totalstress_Profile_unstable.png} &
\includegraphics[width=0.35\textwidth]{Plots/gridStudy/Plots/Totalstress_Profile_stable.png}
\end{tabular}
\caption{Vertical profile of squared friction velocity $u_*^2$ for (a) convective and (b) stable conditions for different resolutions.}
\label{fig:GRAustar2}
\end{figure}

Next, to evaluate the turbulence-resolving capability of the LES, various metrics are considered, including the resolved kinetic energy ($k_{res}$), the subgrid-scale kinetic energy ($k_{sgs}$), the effective resolution index ($\gamma$) \citep{pope2001turbulent}, the one-dimensional energy spectra, and the two-point velocity correlations. These metrics collectively describe how effectively the numerical grid resolves the range of turbulent scales present in the flow.

The resolved kinetic energy is obtained from the trace of the Reynolds stress tensor ($R_{ij} = \overline{u'_i u'_j}$) and defined as $k_{res} = (1/2)(\overline{u'^2} + \overline{v'^2} + \overline{w'^2})$. The subgrid-scale kinetic energy is similarly calculated from the subgrid stress tensor ($\tau_{ij}$) as $k_{sgs} = (1/2)\mathrm{Tr}(\tau_{ij})$. The ratio of resolved to total turbulent kinetic energy defines the effective resolution index $\gamma = k_{res}/(k_{res} + k_{sgs})$, which represents the fraction of turbulence energy directly resolved by the simulation. According to~\citet{pope2001turbulent}, values of $\gamma \ge 0.8$ typically indicate a well-resolved LES. 

However, \citet{wurps2020grid}~and \citet{davidson2009large} noted that effective resolution index $\gamma$ alone may not sufficiently assess grid adequacy, since even under-resolved simulations can produce $\gamma > 0.8$. To obtain a more complete assessment, the one-dimensional energy spectra and two-point correlations of the velocity fluctuations are also examined, providing information on the scale distribution and spatial coherence of turbulence. The one-dimensional streamwise energy spectrum $E(k_x)$ quantifies the distribution of resolved kinetic energy across spatial wavenumbers. It is computed using Fourier transforms of the streamwise velocity fluctuations $u'(x)$ sampled on horizontal planes at $z = 200$~m (i.e., $z/H = 0.1$) and averaged over a two-hour window centered on the analysis time ($\pm 1~\mathrm{hr}$). The smallest and largest resolvable wavenumbers are determined by the domain length and the finest grid spacing, respectively, where $k_{\mathrm{min}} = 1/L_x$ and $k_{\mathrm{max}} = 1/(2\Delta)$. For $L_x = 3~\mathrm{km}$ and $\Delta = 10~\mathrm{m}$, these correspond to $k_{\mathrm{min}} \approx 3.3\times10^{-4}~\mathrm{m^{-1}}$ and $k_{\mathrm{max}} = 5.0\times10^{-2}~\mathrm{m^{-1}}$. On the other hand, spatial coherence is examined through the two-point correlation function of the streamwise velocity fluctuations, defined as $B_{\Delta u}(r) = \langle u'(x)u'(x - r)\rangle / \sigma_u^2$, where $r$ is the separation distance in the streamwise direction and $\sigma_u^2$ is the variance of $u'$. The functional is normalized such that $B_{\Delta u}(0)= 1$ and decreases with increasing $r$, reflecting loss of correlation at larger separations. The correlation can also be derived from the energy spectrum by the Fourier cosine transform, $B_{\Delta u} (r) = (1/\sigma_u^2)\int_0^\infty E(k_x)\cos(2\pi k_x r)\,dk_x$, demonstrating the equivalence between the spatial and spectral representations of turbulence. Following \citet{wurps2020grid}and ~\citet{davidson2009large}, the characteristic turbulent structure size $\sigma$ which one half of the average structure in the flow and it is identified as the distance where $B_{\Delta u} (r) =0.3$. The number of grid points $n_c$ resolving an average turbulent structure is given by $\sigma = n_c \Delta$, where $\Delta$ is the grid spacing.

Figures~\ref{fig:kres}a and~\ref{fig:kres}b presents the vertical profiles of resolved scale turbulent kinetic energy ($k_{res}$) for convective and stable regimes, respectively. Under convective conditions (Fig.~\ref{fig:kres}a), the coarse isotropic grids (30$\times$30$\times$30 and 40$\times$40$\times$40) predict the lowest levels of $k_{res}$ across the entire surface layer of the ABL. In contrast, the grids with finer vertical resolution--10$\times$10$\times$10, 20$\times$20$\times$10, and 30$\times$30$\times$10--show good agreement near the surface, suggesting convergence in that region. Further away from the surface, the 10$\times$10$\times$10 grid captures the highest resolved kinetic energy, followed by 20$\times$20$\times$10 and then 30$\times$30$\times$10. In the stable regime (Fig. \ref{fig:kres}b), $k_{\mathrm{res}}$ appears to be highly sensitive to both grid resolution and mesh aspect ratio. Coarser grids such as 40$\times$40$\times$10 and 30$\times$30$\times$10 produce substantially higher resolved TKE near the surface compared to finer grids like 10$\times$10$\times$10 and 20$\times$20$\times$20, which exhibit convergence with lower $k_{\mathrm{res}}$ values. In principle, the resolved-scale TKE is expected to increase as the grid is refined. However, under stable boundary layer conditions--particularly with meshes having large aspect ratios--higher resolved TKE is observed. This counterintuitive behavior is consistent with earlier findings~\citep{beare2006intercomparison,celik2005, wurps2020grid}, which show that coarse meshes in stable, wall-bounded flows suppress the resolved strain-rate tensor $S_{ij}$, leading to a reduction in resolved dissipation $\varepsilon = 2 \nu_{sgs} S_{ij} S_{ij}$. As a result, less energy is removed by the subgrid-scale model, artificially increasing the resolved TKE.
%
\begin{figure}[htb!]
\centering
\begin{tabular}{cc}
{\large (a)} & {\large (b)} \\
\includegraphics[width=0.35\textwidth]{Plots/gridStudy/Plots/Kres_Profile_unstable.png} &
\includegraphics[width=0.35\textwidth]{Plots/gridStudy/Plots/Kres_Profile_stable.png}
\end{tabular}
\caption{Resolved kinetic energy $k_{res}$ for (a) convective and (b) stable stratification for different resolutions.}
\label{fig:kres}
\end{figure}

Figures~\ref{fig:gamma}a and~\ref{fig:gamma}b illustrate the variation of the effective resolution index $\gamma$ with height for different grid configurations. In the convective regime (Fig.~\ref{fig:gamma}a), the finest grid (10$\times$10$\times$10) reaches $\gamma$ values as high as 0.9, indicating that 90\% of the total turbulent kinetic energy is resolved and satisfies the well-resolved LES criterion proposed by \citet{pope2001turbulent}. Several other grids--specifically 20$\times$20$\times$10, 30$\times$30$\times$10, 20$\times$20$\times$20, and 40$\times$40$\times$10--also achieve $\gamma$ values between 0.8 and 0.9 throughout most of the surface layer, indicating that these resolutions are sufficient for resolving the majority of turbulence under convective conditions. The 30$\times$30$\times$30 grid falls slightly below the 0.8 threshold near the surface but reaches around 0.8 further away from the surface, while the coarsest grid configuration of 40$\times$40$\times$40 fails to achieve $\gamma = 0.8$ at any height, suggesting that it under-resolves turbulence across the entire boundary layer. In contrast, the stable regime (Fig.~\ref{fig:gamma}b) exhibits substantially lower $\gamma$ values across all grids. Even the finest case does not reach the 0.8 threshold, achieving a maximum of just above 0.5. All other grid configurations remain below 0.5, indicating that the flow is under-resolved in terms of turbulent kinetic energy. These results reflect the well-known difficulty of resolving turbulence in stable boundary layers, where the energy-containing eddies are smaller and more intermittent. This under-resolution is consistent with previous findings suggesting that grid spacing in the range of $2$-$5$ m are typically required to capture turbulence in stable boundary layers \citep{wurps2020grid, gadde2021interaction}. The inability of current grid configurations to meet this resolution threshold may explain the overall deficiency in resolving near-surface turbulence under stable stratification.

\begin{figure}[htb!]
\centering
\begin{tabular}{cc}
{\large (a)} & {\large (b)} \\
\includegraphics[width=0.35\textwidth]{Plots/gridStudy/Plots/Gamma_Profile_unstable.png} &
\includegraphics[width=0.35\textwidth]{Plots/gridStudy/Plots/Gamma_Profile_stable.png}
\end{tabular}
\caption{Effective resolution index $\gamma$ for (a) convective and (b) stable regimes under different grid resolutions. The two vertical black dashed lines represent $80\%$ and $90\%$ $\gamma$ respectively. }
\label{fig:gamma}
\end{figure}

Figures~\ref{fig:spectra}a and ~\ref{fig:spectra}b  shows the normalized one-dimensional energy spectra $E(k_x)/E(0)$ as a function of the streamwise wavenumber $k_x$~($\mathrm{m}^{-1}$) for both convective and stable regimes. In the convective regime (Fig.~\ref{fig:spectra}a), the finest grid (10$\times$10$\times$10) exhibits the widest inertial subrange, indicating that a broad range of turbulent structures is being resolved. The 20$\times$20$\times$10 grid follows closely, capturing a slightly narrower inertial range. The grids 30$\times$30$\times$10, 30$\times$30$\times$30, and 40$\times$40$\times$10 show similar trends, with more limited inertial scaling and earlier spectral roll-off, suggesting intermediate resolution quality. The coarsest grid (40$\times$40$\times$40) captures the least amount of the inertial subrange, highlighting its inability to resolve multiscale turbulence effectively. In the stable regime (Fig.~\ref{fig:spectra}b), the differences among grids are more pronounced. The coarsest grid (40$\times$40$\times$40) fails to capture any identifiable inertial subrange. All other grids capture only a small portion of it, with the finest grid (10$\times$10$\times$10) exhibiting the clearest signature of limited inertial scaling. This is followed by the 20$\times$20$\times$10 grid, which also retains part of the inertial range but with reduced energy content. These observations are consistent with the previous $\gamma$ analysis, where only the finer grids were able to resolve a substantial fraction of the turbulent kinetic energy in the stable boundary layer.

\begin{figure}[htb!]
\centering
\begin{tabular}{cc}
{\large (a)} & {\large (b)} \\
\includegraphics[width=0.48\textwidth]{Plots/gridStudy/Plots/old/spectrum_unstable.png} &
\includegraphics[width=0.48\textwidth]{Plots/gridStudy/Plots/old/spectrum_stable.png}
\end{tabular}
\caption{Streamwise velocity spectra $E(k_x)$ for (a) convective and (b) stable regimes.}
\label{fig:spectra}
\end{figure}

Figures~\ref{fig:corr}a and~\ref{fig:corr}b displays the two-point correlation functions \( B_{\Delta u} (r) \) for convective and stable atmospheric conditions, computed at \( z/H = 0.1 \) and averaged over a two-hour window. In the convective regime (Fig.~\ref{fig:corr}a), all grids exhibit a gradual decay of correlation with streamwise separation, consistent with active turbulence. Although the decay trends are broadly similar, finer grids such as 10$\times$10$\times$10 exhibit slightly faster decorrelation, followed closely by 20$\times$20$\times$10, indicating their ability to capture smaller coherent structures. The location where \( B_{\Delta u} (r) = 0.3 \), marked by the dashed black line, represents the typical structure size. Differences among grids at this threshold are subtle in the convective case, suggesting that all grids resolve the dominant structures reasonably well. In contrast, Fig.~\ref{fig:corr}(b) for the stable regime shows a more pronounced sensitivity to grid resolution. Finer grids--specifically 10$\times$10$\times$10 and 20$\times$20$\times$10--exhibit significantly steeper decay curves and reach \( B_{\Delta u} (r) = 0.3 \) much earlier than coarser grids. This indicates a finer resolution of the smaller and weaker turbulent structures that persist under stable conditions. Coarse grids such as 30$\times$30$\times$30 and 40$\times$40$\times$40 decay more gradually and fail to resolve the typical turbulent length scales, resulting in lower values of \( n_c \), the number of grid points per resolved structure. These observations reinforce the importance of adequate grid resolution when simulating stable boundary layers, where turbulence is more intermittent and spatially compact. It is worth noting that if the computational domain is too small, periodic boundary conditions can artificially prevent \( B_{\Delta u} (r) \) from decaying to zero at large separations. However, our domain size is sufficiently large to avoid such effects in the reported results.

\begin{figure}[htbp]
\centering
\begin{tabular}{cc}
{\large (a)} & {\large (b)} \\
\includegraphics[width=0.45\textwidth]{Plots/gridStudy/corr_unstable.pdf} &
\includegraphics[width=0.45\textwidth]{Plots/gridStudy/corr_stable.pdf}
\end{tabular}
\caption{Two-point correlation functions \( B_{\Delta u} \) for (a) convective and (b) stable regimes at \( z/h = 0.1 \).}
\label{fig:corr}
\end{figure}

Finally, the analysis of grid resolution based on turbulence resolving metrics and two point correlation functions indicates that the $20\times20\times10$ grid performs comparably to the finer $10\times10\times10$ grid under convective (unstable) boundary layer conditions. This grid captures more than 80\% of the resolved kinetic energy, as demonstrated by the LES effective resolution index ($\gamma$), and accurately reproduces the sub inertial range of the turbulence energy spectra. In stable regimes, the guideline $\gamma \ge 0.8$ is not fully met, since none of the tested resolutions reach this threshold. However, the energy spectra and two point correlation functions for the $20\times20\times10$ grid remain in reasonable agreement with those obtained on the finer grid, indicating that the dominant turbulent structures and their characteristic length scales are well represented. On this basis, the $20\times20\times10$ grid is chosen for the present complete diurnal cycle simulations. Several factors motivate this choice. First, performing complete diurnal cycle simulations at a finer resolution (for example $10\times10\times10$) is computationally prohibitive for multiple cases of the various profile assimilation methods. Second, for wind energy and mesoscale to microscale coupling studies, the key quantities influencing power generation, such as the mean wind profile, shear stress, and momentum flux, are well represented at this resolution, while the resolved kinetic energy primarily affects secondary quantities such as dynamic loading. 

Taken together, these findings indicate that the $20\times20\times10$ grid provides a reasonable compromise between physical fidelity and computational cost. It offers sufficient accuracy to reproduce the dominant turbulent structures and key mean flow features across the different stability regimes, while remaining practical for extended diurnal cycle simulations. This grid is therefore well suited for simulations over larger domain extents, especially when the primary interest lies in the large scale dynamics of the atmospheric boundary layer.

\bibliographystyle{spbasic_updated}     
\bibliography{ref} 

\end{document}
